\chapter{Symmetry Properties of the Fourier Transform}
\label{ch:c4}

Fourier analysis is intimately connected with the action of the group of translations on Euclidean space. If we were to restrict our study only to questions related to this translation structure we would present little that cannot be found in a treatment of abstract harmonic analysis on locally compact Abelian groups and, at the same time, lose much of its elegant generality. Harmonic analysis in Euclidean spaces, however, is richer because of its connection with several classes of transformations: the dilations and the rotations as well as the translations. This connection has already appeared implicitly in our discussion of harmonic and holomorphic functions. We now wish to make quite explicit the further connection between the Fourier transform, dilations, and rotations.

The starting point for this is the observation that the Fourier transform has a very simple transformation law under dilations and that, in addition, it commutes with the action of rotations. This second property leads to a direct sum decomposition of $L^{2}(E_n)$ into various parts, each transforming in a specific manner with respect to the action of rotations. The Fourier transform preserves this direct sum and its restriction to each of these parts can be identified with a classical Bessel transform. This decomposition is obtained by a passage to polar coordinates enabling us first to decompose the $L^{2}$ spaces of functions on the surface of the unit sphere of $E_n$. This decomposition is obtained in terms of the restriction to the surface of the unit sphere of homogeneous harmonic polynomials. The study of these functions, the spherical harmonics, is very natural in view of the considerations mentioned above because the Laplacian is both translation and rotation invariant while the homogeneity of these polynomials reflects the action of dilations.

This chapter is organized in the following way: the first section treats the special two-dimensional case, which is particularly simple since the decomposition of $L^{2}$ of the surface of the unit sphere is obtained in terms of the usual one-dimensional Fourier series expansions of square integrable functions. The second section generalizes these expansions to n dimensions by developing the theory of spherical harmonics. The third section begins with the study of the space of radial functions, the simplest of the summands occurring in the above mentioned decomposition of $L^{2}(E_{n})$, and then treats the rest of the direct sum decomposition. The fourth section deals with some applications of

this theory involving several important identities in the theory of potentials and singular integrals, where the interplay of dilations and rotations\index{rotation@rotation} is evident.

\section{The Decomposition of  $L^{2}(E_{2})$ into Subspaces Invariant Under the Fourier Transform}

Our principal goal in this chapter is to study in much greater detail the action of the Fourier transform on functions defined on $E_n$. This will be done by finding a natural direct sum decomposition of $L^2(E_n)$ that is preserved by the Fourier transform and, then, by making a detailed analysis of how the latter acts on each of the summands. In this connection we will be led to a development of the basic properties of Bessel functions and spherical harmonics. These are of fundamental importance, not only in other aspects of harmonic analysis, but in many other branches of analysis as well.

In one dimension we have a simple natural decomposition of $L^{2}(E_{1})$. The subspace of even functions and the subspace of odd functions are easily seen to be invariant under the Fourier transform, they are orthogonal\index{orthogonal transformation@orthogonal transformation} to each other and their direct sum is all of $L^{2}(E_{1})$. These properties, obtained by decomposing a function $f$ into the sum $f = f_{e} + f_{0}$, where the $even\ part\ f_{e}$ is defined by $f_{e}(x) = [f(x) + f(-x)]/2$ and the $odd\ part\ f_{0}$ by $f_{0}(x) = [f(x) - f(-x)]/2$, are clearly also valid in $n$ dimensions. A much more interesting, though more complicated, situation arises when we consider the following higher dimensional generalization of the even part of a function. Given a locally integrable function $f$ on $E_{n}$, its $radial\ part$ is the function $\varphi$ defined by\index{radial part of a function@radial part of a function}

\[
\varphi(x)=\frac{1}{\omega_{n-1}}\intop_{\Sigma_{n-1}}f(r x^{\prime})d x^{\prime},
\]

where, we recall, $r = |x|, x' = x/r$ (when $x \neq 0$) and $\omega_{n-1}$ is the area of the surface of the unit sphere $\Sigma_{n-1}$ of $E_n$. It is clear that $\varphi$ is a radial function; that is, $\varphi$ depends only on $r = |x|$. When $n = 1$ we see that $\varphi = f_e$.

This notion, which is certainly natural when considering polar coordinates, is also appropriate when studying the action of the Fourier transform. The following remarks, from which we will deduce that the Fourier transform of a radial function is radial, will make this clear. We first observe that a function $f$, defined on $E_n$, is radial if and only if $f(\rho x) = f(x)$ for all orthogonal transformations\footnote{We recall that $\rho$ is orthogonal if it is a linear operator on $E_n$ that preserves the inner product: $\rho x \cdot \rho y = x \cdot y$ for all $x, y$ in $E_n$. If $\det \rho = 1$, $\rho$ is called a rotation. It is also true that $f$ is radial if and only if $f(\rho x) = f(x)$ for all rotations $\rho$ and all $x$ in $E_n$, when $n > 1$.} $\rho$ of $E_n$ and all $x$ in $E_n$. The basic property

of the Fourier transform with respect to orthogonal transformations is the following result:

\begin{theorem}
\label{thm:c4-1-1}
The Fourier transform, $\mathcal{F}$, commutes with orthogonal transformations. That is, if $\rho$ is an orthogonal transformation, let $R_\rho$ be the mapping taking $a$ function $f$ on $E_n$ into $a$ a function $g$ whose values are $g(x) = (R_\rho f)(x) = f(\rho x)$ for $x \in E_n$, then, whenever $f \in L^1(E_n)$,
\[
\hat{g}(t)=(\mathcal{F}g)(t)=(\mathcal{F}R_{\rho}f)(t)=(R_{\rho}\mathcal{F}f)(t)=(\mathcal{F}f)(\rho t)=\hat{f}(\rho t);
\]

that is, the operators $\mathscr{F}$ and $R_{\rho}$ commute: $\mathscr{F}R_{\rho} = R_{\rho}\mathscr{F}$.
\end{theorem}

\textsc{Proof.}\ Since the adjoint of $\rho$ is also its inverse, and the Jacobian in the change of variables $w = \rho x$ is one, we have

\[
\begin{aligned}{\hat{g}(t)}&{{}=\intop_{E_{n}}e^{-2\pi i t\cdot x}f(\rho x)\:d x=\intop_{E_{n}}e^{-2\pi i t\cdot\rho^{-1}w}f(w)\:d w}\\ {}&{{}=\intop_{E_{n}}e^{-2\pi i\rho t\cdot w}f(w)\:d w=\hat{f}(\rho t).}\\ \end{aligned}
\]

Since it is true that whenever $|x_1| = |x_2|$ for two points of $E_n$ there is an orthogonal transformation $\rho$ such that $\rho x_1 = x_2$, we immediately obtain the above mentioned property of the Fourier transform:

\begin{corollary}
\label{cor:c4-1-2}
If $f$ is a radial function in $L^{1}(E_{n})$ then $\hat{f}$ is also radial.
\end{corollary}

This result extends in an obvious way to $L^2(E_n)$. We thus obtain a direct sum decomposition of this space which is left invariant by the Fourier transform simply by considering the subspace $\mathfrak{S}^0$ of those functions equal almost everywhere to radial functions, and the orthogonal complement of $\mathfrak{S}^0$ (this is an immediate consequence of the Plancherel theorem). Consequently, it is natural to ask whether there exists a simple description of the space of functions orthogonal to the radial ones and, also, if we can obtain more detailed information on the action of the Fourier transform on this subspace. We first consider these questions in the two dimensional case.

Let us choose an $f$ in $L^2(E_2)$ and let us use the standard identification of the points $(x, y)$ in $E_2$ with the complex numbers $x + iy = z = re^{i\theta}$. Fubini's theorem implies that, for almost every $r$, $f(re^{i\theta})$ defines a square integrable function of $\theta$. Consequently, we have a Fourier series expansion

\begin{equation}\label{eq:c4-1-3}
\begin{array}{c} f(re^{i\theta}) \sim \displaystyle \sum_{k = - \infty}^{\infty} f_{k}(r)e^{ik\theta} \end{array}
\end{equation}

which converges to $f(re^{i\theta})$ in the $L^2$-norm for almost every $r$. Moreover, $\sum_{-\infty}^{\infty}|f_k(r)|^2 = (2\pi)^{-1}\int_0^{2\pi}|f(re^{i\theta})|^2\,d\theta$ a.e.; hence, by the Lebesgue

monotone convergence theorem.

\[
\lim_{n\to\infty}\int\limits_{0}^{\infty}\sum_{-n}^{n}|f_{k}(r)|^{2}r dr=\frac{1}{2\pi}\int\limits_{0}^{\infty}\left\{\int\limits_{0}^{2\pi}|f(re^{i\theta})|^{2}d\theta\right\}r dr=\frac{1}{2\pi}\|f\|_{2}^{2}.
\]

Putting $g_k(z) = f_k(r)e^{ik\theta}$, $k = 0, \pm1, \pm2, \ldots$, we thus obtain

\[
\begin{aligned}\lim_{n\to\infty}\int\limits_{E_{2}}|f(z)-\sum_{-n}^{n}g_{k}(z)|^{2}dz\\=\lim_{n\to\infty}2\pi\int\limits_{0}^{\infty}\left[\frac{1}{2\pi}\left\{\int\limits_{0}^{2\pi}|f(re^{i\theta})|^{2}d\theta\right\}-\sum_{-n}^{n}|f_{k}(r)|^{2}\right]r dr=0.\end{aligned}
\]

This, together with the orthogonality relations satisfied by the exponential functions $e^{ik\theta}, k = 0, \pm1, \pm2, \ldots$, shows that we have the direct sum decomposition

\begin{equation}\label{eq:c4-1-4}
L^{2}(E_{2})=\sum_{k=-\infty}^{\infty}\oplus\mathfrak{H}^{k},\footnote{By such a direct sum we mean that each of the subspaces $\mathfrak{H}^{k}$ are closed in $L^{2}(E_{2})$, are mutually orthogonal and the closure of their linear span consists of all of $L^{2}(E_{2})$.}
\end{equation}

where $\mathfrak{J}^k = \{g \in L^2(E_2); g(z) = f(r)e^{ik\theta}\}$ a.e. for some measurable function satisfying $\int_0^\infty |f(r)|^2 dr < \infty$. This notation is consistent with that used above since $\mathfrak{J}^0$, in either case, is precisely the subspace of those square integrable functions that are equal almost everywhere to radial functions. Moreover, we see from \eqref{eq:c4-1-4} that the orthogonal complement of $\mathfrak{J}^0$ is the direct sum of the subspaces $\mathfrak{J}^k, k \neq 0$.

We know that the Fourier transform maps $\tilde{\mathfrak{H}}^{0}$ into itself. It is easy to see that it also maps each of the other spaces $\mathfrak{H}^{k}, k = \pm 1, \pm 2, \ldots$, into themselves: First suppose $g$ belongs to $\tilde{\mathfrak{H}}^{k} \cap L^{1}(E_{2})$ and set $h(z) = g(e^{i\phi}z)$ for $\phi$ fixed. We then have $h(z) = e^{ik\phi}g(z)$ for almost every $z$. On the other hand, multiplication by $e^{i\phi}$ is a rotation of $E_{2}$ and, since the Fourier transform commutes with rotations (Theorem \ref{thm:c4-1-1}), we have

\[
\hat{g}(e^{i\phi}w)=\hat{h}(w)=e^{i k\phi}\hat{g}(w)
\]

for all $w$ and $\phi$. Putting $w = r \geq 0$ we see that $\hat{g}$ also belongs to $\mathfrak{H}^k$. Since $\mathfrak{H}^k \cap L^1(E_2)$ is dense in the closed subspace $\mathfrak{H}^k$ this shows that all of this last space is mapped into itself by the Fourier transform. Because of the Plancherel theorem we can deduce that all the spaces $\mathfrak{H}^k$ are mapped onto each other.

Therefore, we have found a direct sum decomposition of $L^{2}(E_{2})$, namely \eqref{eq:c4-1-4}, that is left invariant by the Fourier transform. This answers our first

question. In order\index{Bessel function@Bessel function!of integral order@of integral order} to answer the second one we now engage in a more detailed study of how the Fourier transform acts on each of the summands $\mathfrak{H}^{k}$.

Let us choose an $f$ belonging to $\mathfrak{S}^k$; this $f$ must have the form $f(z) = f_0(r)e^{ik\theta}$ for almost every $z = re^{i\theta}$. Since $\hat{f}$ must also belong to $\mathfrak{S}^k$ it has a similar form; that is, $\hat{f}(w) = F_0(R)e^{ik\phi}$ for almost every $w = Re^{i\phi}$. Assuming $f$ is integrable we shall now compute $F_0$ explicitly (our results extend in an obvious way to square integrable $f$). Putting $w = Re^{i0} = R$, we have

\[
\begin{aligned}{F_{0}(R)}&{{}=\hat{f}(Re^{i0})=\int_{0}^{\infty}f_{0}(r)\Biggl\{\int_{0}^{2\pi}e^{-2\pi i R r\cos\theta}e^{i k\theta}d\theta\Biggr\}r\;d r}\\ {}&{{}=(-i)^{k}2\pi\int_{0}^{\infty}f_{0}(r)\Biggl\{\frac{1}{2\pi}\int_{0}^{2\pi}e^{2\pi i R r\sin\theta}e^{-i k\theta}d\theta\Biggr\}r\;d r.}\\ \end{aligned}
\]

Thus, we see that an explicit expression for $F_0(R)$ in terms of $f_0$ involves the Fourier coefficients of the function of $\theta$ given by the expression $e^{i t \sin \theta}$. These coefficients, which depend on t, are known as the Bessel functions $J_k(t)$. That is, the expression

\[
J_{k}(t) = \frac{1}{2\pi} \int\limits_{0}^{2\pi}e^{i t\sin\theta}e^{-i k\theta} d\theta
\]

defines the Bessel function $J_k$, $k = 0$, $\pm 1$, $\pm 2$, $\ldots$. By a simple change of variables we obtain the relation, valid for all integers $k$,

\begin{equation}\label{eq:c4-1-5}
J_{k}(t)=(-1)^{k}J_{-k}(t).
\end{equation}

Hence, we can state our main result on the action of the Fourier transform on the spaces $\mathfrak{H}^k$ in the following way:

\begin{theorem}
\label{thm:c4-1-6}
Suppose $f$ belongs to $L^{1}(E_{2})$ and $f(z)=f_{0}(r)e^{ik\theta}$, where $z=re^{i\theta}$; then $\hat{f}(w)=F_{0}(R)e^{ik\phi}$, where $w=Re^{i\phi}$ and $F_{0}(R)=2\pi i^{k}\int_{0}^{\infty}f_{0}(r)J_{-k}(2\pi Rr)r dr=2\pi(-i)^{k}\int_{0}^{\infty}f_{0}(r)J_{k}(2\pi Rr)r dr$.
\end{theorem}

\section{Spherical Harmonics}

In order to extend this decomposition to $n$ dimensions it would be desirable to be able to represent functions defined on the surface of the unit sphere by expansions similar to Fourier series. This would enable us to obtain a development for $f(x) = f(rx')$ analogous to \eqref{eq:c4-1-3}. We shall show that this can be done by letting the rôle of the exponentials $e^{ik\theta}, k = 0, \pm 1, \pm 2, \ldots$, be assumed by a class of functions called the spherical

harmonics. Before defining these functions we make some simple observations about trigonometric series $s(\theta) = \sum_{k=-\infty}^{\infty} c_k e^{ik\theta}$.

We consider such a series to be the sequence of symmetric partial sums $s_n(\theta) = \sum_{k=-n}^{n} c_k e^{ik\theta} = c_0 + \sum_{k=1}^{n} (c_k e^{ik\theta} + c_{-k} e^{-ik\theta})$. These sums are real valued provided $c_k = \tilde{c}_{-k}$. In this case $s_n(\theta)$ is the restriction to the unit circle $\Sigma_1$ of the real part $u_n(z)$ of the polynomial $p_n(z) = c_0 + 2\sum_{k=1}^{n} c_k z^k$, $z = re^{i\theta} \in C_1$. Since there is no restriction imposed on the complex coefficients $c_1, c_2, \ldots, c_n$ and $c_0$ can be any real number, it follows that $u_n(z) = u_n(x + iy)$ is an arbitrary real harmonic polynomial of degree $n$ in the variables $x$ and $y$. The summands $c_0$, $\text{Re}\{c_k z^{k}\} = q_k(z)$, $k=1,2,\ldots$, are then the general homogeneous\footnote{A function $f$ defined on $E_n$ is said to be homogeneous of degree $k$ if $f(ax) = a^k f(x)$ for all $x \in E_n$ and $a > 0$.} real harmonic polynomials of degree $k$. Thus, the linear combinations $Y^{(k)}(e^{i\theta}) = c_k e^{ik\theta} + \tilde{c}_k e^{-ik\theta} = a_k \cos k\theta + b_k \sin k\theta$, with $c_k = (a_k - ib_k)/2$, are precisely the restrictions to the unit circle $|z| = 1$ of such polynomials $q_k(z) = q_k(re^{i\theta})$. In particular, since each real valued $L^2$-function on the unit circle can be represented (in the $L^2$ norm) by its Fourier series, whose terms are such restrictions, we see that $L^2(\Sigma_1)$ is the closure of the span of these functions $Y^{(k)}, k=0,1,2,\ldots$

In general, the restriction to the surface of the unit sphere $\Sigma_{n-1}$ of a homogeneous harmonic polynomial of degree $k$ is called a spherical harmonic\index{spherical harmonic of degree k@spherical harmonic of degree $k$} of degree $k$. This section will be devoted to the development of important properties of spherical harmonics. This will enable us, in particular, to show how $L^2(E_n)$ can be decomposed into a direct sum analogous to \eqref{eq:c4-1-4}, where the summands $\mathfrak{G}_k$ consist of the subspaces of $L^2(E_n)$ that are spanned by functions that are products of radial functions times spherical harmonics of degree $k$.\footnote{When $n = 2$ the space $\mathfrak{H}_{k}$ is spanned by $\mathfrak{H}^{k}$ and $\mathfrak{H}^{-k}, k = 0, 1, 2, \ldots$.} The Fourier transform leaves these spaces $\mathfrak{G}_k$ invariant and we shall describe its behavior on them in detail, thus obtaining an extension of Theorem \ref{thm:c4-1-6}.

We let $\mathcal{P}_k$ be the set of all homogeneous polynomials of degree $k$ on $E_n$ with complex coefficients. Thus, if $P \in \mathcal{P}_k$ then

\[
P(x)=\sum_{|\alpha|=k}c_\alpha x^\alpha,
\]

where (as in Chapter I) $\alpha$ denotes an $n$-tuple $(\alpha_1, \alpha_2, \ldots, \alpha_n)$ of nonnegative integers, $|\alpha| = \alpha_1 + \alpha_2 + \cdots + \alpha_n$ and $x^\alpha = x_{1}^{\alpha_1} x_{2}^{\alpha_2} \cdots x_{n}^{\alpha_n}$. It is clear that the set of monomials $x^\alpha$, $|\alpha| = k$, is a basis for this space. The number of such monomials, on the other hand, is precisely the number, $d_k$, of ways an $n$-tuple $\alpha = (\alpha_1, \alpha_2, \ldots, \alpha_n)$ of nonnegative integers can be chosen so

that $\alpha_1 + \alpha_2 + \cdots + \alpha_n = k$. It is not hard to determine $d_k$. Let us select $n - 1$ boxes out of a linearly ordered array of $n + k - 1$ boxes and let us fill each of the remaining $k$ boxes with a ball. Then, there will be $\alpha_1$ balls preceding the first box we chose, $\alpha_2$ balls in the boxes between the first and second one we chose and so on until we reach the $\alpha_n$ balls that follow the last box we chose. In this way we obtain $n$ nonnegative integers $\alpha_1, \alpha_2, \ldots, \alpha_n$ satisfying $\alpha_1 + \alpha_2 + \cdots + \alpha_n = k$ and all such $n$-tuples are obtainable in this manner. Thus, there are precisely as many such $n$-tuples as there are ways of selecting $n - 1$ boxes out of a collection of $n + k - 1$. Consequently, the dimension of $\mathcal{P}_k$ is

\[
d_{k}=\binom{n+k-1}{n-1}=\binom{n+k-1}{k}=\frac{(n+k-1)!}{(n-1)!k!}.
\]

We introduce an inner product $\langle P, Q \rangle$ on $\mathcal{P}_k$ by letting $\langle P, Q \rangle = P(D) \overline{Q}$ for all $P, Q$ in $\mathcal{P}_k$, where $P(D)$ is the differential operator introduced in Chapter I (see \eqref{eq:c1-1-9}). Since $P$ and $Q$ are homogeneous polynomials of the same degree, $\langle P, Q \rangle$ is scalar valued; furthermore it is clearly linear in the first variable, conjugate linear in the second variable and hermitian symmetric. In order to verify that we, indeed, have an inner product, we must, therefore, show that $\langle P, P \rangle \geq 0$, with equality only if $P = 0$. But if $\alpha = (\alpha_1, \ldots, \alpha_n) \neq (\beta_1, \ldots, \beta_n) = \beta$ then

\[
\biggl(\frac{\partial^{\alpha_{1}}}{\partial x_{1}^{\alpha_{1}}}\frac{\partial^{\alpha_{2}}}{\partial x_{2}^{\alpha_{2}}}\cdots\frac{\partial^{\alpha_{n}}}{\partial x_{n}^{\alpha_{n}}}\biggr)x_{1}^{\beta_{1}}x_{2}^{\beta_{2}}\cdots x_{n}^{\beta_{n}}=0,
\]

while this derivative equals $\alpha_1! \alpha_2! \cdots \alpha_n! = \alpha!$ when $\alpha = \beta$. Consequently, if $P(x) = \sum_{|\alpha| = k} c_\alpha x^\alpha$ then $\langle P, P \rangle = \sum_{|\alpha| = k} |c_\alpha|^2 \alpha!$. But this last expression vanishes if and only if all the coefficients $c_\alpha$ are 0.

We shall use this inner product in order to establish the following basic result:

\begin{theorem}
\label{thm:c4-2-1}
If $P \in \mathcal{P}_k$ then
\[
P(x)=P_{0}(x)+\left|x\right|^{2}P_{1}(x)+\cdots+\left|x\right|^{2l}P_{l}(x),
\]

where $P_j$ is a homogeneous harmonic polynomial of degree k - 2j, j = 0, 1,..., l.
\end{theorem}

\textsc{Proof.}\ Any polynomial of degree less than 2 is harmonic; thus, we may assume that $k \geq 2$. Consider the linear mapping $\varphi: \mathcal{P}_k \to \mathcal{P}_{k-2}$ defined by letting $\varphi(P) = \Delta P$ for $P \in \mathcal{P}_k$, where $\Delta$ is the Laplace operator. We first show that $\varphi$ maps $\mathcal{P}_k$ onto $\mathcal{P}_{k-2}$. If this were not the case we could find a nonzero $Q \in \mathcal{P}_{k-2}$ that is orthogonal to $\mathcal{R}(\varphi)$, the range of $\varphi$. That is

\[
\overline{\langle\Delta P,\overline{Q}\rangle}=\langle Q,\Delta P\rangle=0
\]

for all $P \in \mathcal{P}_k$. In particular, this must be true for $P(x) = |x|^2 Q(x)$. Thus

\[
0=\langle Q,\Delta P\rangle=Q(D)\overline{{\Delta P}}=\Delta Q(D)\overline{{P}}=P(D)\overline{{P}}=\langle P,P\rangle.
\]

But this is impossible since $P \neq 0$.

Let $\mathcal{A}_j \subset \mathcal{P}_j, j \geq 2$, be the class of all polynomials in $\mathcal{P}_j$ that are harmonic. We claim that $\mathcal{P}_j$ is the orthogonal direct sum of $\mathcal{A}_j$ and $\mathcal{B}_j = |x|^2 \mathcal{P}_{j-2} = \{P(x) \in \mathcal{P}_j; P(x) = |x|^2 Q(x) \quad \text{with} \quad Q \in \mathcal{P}_{j-2}\}$. If $R(x) = |x|^2 Q(x)$ with $Q \in \mathcal{P}_{j-2}$, then $\langle R, P \rangle = 0$ for all $Q \in \mathcal{P}_{j-2}$ if and only if $Q(D) \Delta \bar{P} = 0$ for all $Q \in \mathcal{P}_{j-2}$, which is true if and only if $\langle Q, \Delta P \rangle = 0$ for all $Q \in \mathcal{P}_{j-2}$, which in turn is true if and only if $\Delta P = 0$.

In particular, for $j = k$ and $P \in \mathcal{P}_k$ we have $P(x) = P_0(x) + |x|^2 Q(x)$ with $P_0$ harmonic and $Q \in \mathcal{P}_{k-2}$. Applying our result again for $j = k - 2$ we obtain the decomposition $Q(x) = P_1(x) + |x|^2 Q_1(x)$ with $P_1$ harmonic and $Q_1$ in $\mathcal{P}_{k-4}$. Thus, $P(x) = P_0(x) + |x|^2 P_1(x) + |x|^4 Q_1(x)$ and it is clear that the theorem follows by induction.

An immediate consequence of Theorem \ref{thm:c4-2-1} is:

\begin{corollary}
\label{cor:c4-2-2}
The restriction to the surface of the unit sphere $\Sigma_{n-1}$ of any polynomial of $n$-variables is a sum of restrictions to $\Sigma_{n-1}$ of harmonic polynomials.\footnote{The reader should note that, when n = 2, this fact follows easily from the observations we made at the beginning of this section.}
We let $\mathcal{H}_k$ denote the space of spherical harmonics of degree $k$. $\mathcal{H}_k$ coincides with the collection of all restrictions to $\Sigma_{n-1}$ of the members of $\mathcal{A}_k$. If we consider the restriction $Y$ of $P$ in $\mathcal{A}_k$ (for $x' \in \Sigma_{n-1}$, $Y(x') = P(x')$) then, because of the homogeneity of $P$, $P(x) = |x|^k Y(x/|x|)$ for $x \neq 0$. Thus, the restriction map $P \to Y$ has a trivial kernel and must, therefore, be an isomorphism of $\mathcal{A}_k$ onto $\mathcal{H}_k$. In particular,
\[
\begin{aligned}\dim\mathcal{H}_{k}&=\dim\mathcal{A}_{k}=\dim\mathcal{P}_{k}-\dim\mathcal{P}_{k-2}=d_{k}-d_{k-2}\\&=\binom{n+k-1}{k}-\binom{n+k-3}{k-2}\end{aligned}
\]

for $k \geq 2$. Moreover, $\dim \mathcal{H}_0 = d_0 = 1$ and $\dim \mathcal{H}_1 = d_1 = n$. At the beginning of this section we found that when $n = 2$ the space $\mathcal{H}_k$ is spanned by the two functions $\cos k\theta$ and $\sin k\theta$; thus $\dim \mathcal{H}_k = 2$ for all $k \geq 1$. This is consistent with the result just obtained:
\[
\begin{pmatrix}2+k-1\\k\end{pmatrix}-\begin{pmatrix}k-1\\k-2\end{pmatrix}=(k+1)-(k-1)=2.
\]

\end{corollary}

When n = 3 we obtain $\dim \mathcal{H}_k = 2k + 1$ for $k \geq 0$.

The space $\mathcal{A}_k$ is called the space of solid spherical\index{spherical harmonic@spherical harmonic!solid@solid} harmonics. Sometimes, in order to emphasize the distinction between $\mathcal{H}_k$ and $\mathcal{A}_k$, the members of $\mathcal{H}_k$ are referred to as the surface spherical\index{spherical harmonic@spherical harmonic!surface@surface} harmonics.

\begin{corollary}
\label{cor:c4-2-3}
The collection of all finite linear combinations of elements of $\bigcup_{k=0}^{\infty} \mathcal{H}_{k}$ is
(i) dense in the space of all continuous functions on $\Sigma_{n-1}$ with respect to the $L^\infty$ norm;
(ii) dense in $L^{2}(\Sigma_{n-1})$.
\end{corollary}

\textsc{Proof.}\ Since the space of continuous functions is dense in $L^2(\Sigma_{n-1})$ it is easy to see that (i) implies (ii): given $\varepsilon > 0$ and $f \in L^2(\Sigma_{n-1})$ choose a continuous $g$ such that $\|f - g\|_2 < \varepsilon/2$. If (i) is true we can then find a finite linear combination, $h$, of elements of $\bigcup_{k=0}^{\infty} \mathcal{H}_k$ such that $\|g - h\|_\infty < \varepsilon/(2\sqrt{\omega_{n-1}})$. Thus, $\|f - h\|_2 \leq \|f - g\|_2 + \|g - h\|_2 < \varepsilon/2 + \sqrt{\omega_{n-1}} \|g - h\|_\infty < \varepsilon/2 + \varepsilon/2 = \varepsilon$.

On the other hand, (i) follows from the Weierstrass approximation theorem: If $g$ is continuous on $\Sigma_{n-1}$ we can approximate it uniformly by polynomials restricted to $\Sigma_{n-1}$. But, by \eqref{cor:c4-2-2}, these restrictions are finite linear combinations of elements of $\bigcup_{k=0}^{\infty} \mathcal{H}_k$.

\begin{corollary}
\label{cor:c4-2-4}
If $Y^{(k)}$ and $Y^{(l)}$ are spherical harmonics of degrees k and l, with $k \neq l$, then
\[
\int\limits_{\Sigma_{n-1}}Y^{(k)}(x^{\prime})Y^{(l)}(x^{\prime})d x^{\prime}=0.
\]

\end{corollary}

\textsc{Proof.}\ For $x \neq 0$ in $E_n$ let $r = |x|$ and $x' = x/r$. We then define $u(x) = r^k Y^{(k)}(x')$ and $v(x) = r^l Y^{(l)}(x')$ for $x \neq 0$ and we put $u(0) = 0 = v(0)$ if neither $k$ nor $l$ is 0, while if, say, $k = 0$ then $Y^{(k)}$ is a constant function and we let $u(0)$ be this constant. The derivatives of $u$ and $v$ at $x' \in \Sigma_{n-1}$ in the direction of the outward directed normal to $\Sigma_{n-1}$ are $(dr^k/dr) Y^{(k)}(x') = k Y^{(k)}(x')$ and $(dr^l/dr) Y^{(l)}(x') = l Y^{(l)}(x')$. Moreover, since $Y^{(k)}$ and $Y^{(l)}$ are (surface) spherical harmonics, $u$ and $v$ are solid spherical harmonics. Thus, by Green's theorem,

\[
\begin{aligned}&0=\int\limits_{|x|\leq1}(u\;\Delta v-v\;\Delta u)dx\\&=\int\limits_{\Sigma_{n-1}}\left(u\frac{\partial v}{\partial r}-v\frac{\partial u}{\partial r}\right)dx^{\prime}\\&=\int\limits_{\Sigma_{n-1}}\left(lY^{(k)}(x^{\prime})Y^{(l)}(x^{\prime})-k\;Y^{(k)}(x^{\prime})Y^{(l)}(x^{\prime})\right)dx^{\prime}\\&=\left(l-k\right)\int\limits_{\Sigma_{n-1}}Y^{(k)}(x^{\prime})Y^{(l)}(x^{\prime})dx^{\prime}.\end{aligned}
\]

The result now follows by dividing by $(l - k)$, which we can do since $l \neq k$.

Let us consider $\mathcal{H}_k$ as a subspace of $L^2(\Sigma_{n-1})$ with the inner product $(f, g) = \int_{\Sigma_{n-1}} f(x') \overline{g(x')} \, dx'$. If $\{Y_1^{(k)}, \ldots, Y_{a_k}^{(k)}\}$, $a_k = d_k - d_{k-2}$, is an orthonormal basis of $\mathcal{H}_k$ then it follows from Corollary \ref{cor:c4-2-4} that the collection $\bigcup_{k=0}^{\infty} \{Y_1^{(k)}, \ldots, Y_{a_k}^{(k)}\}$ is an orthonormal basis of $L^2(\Sigma_{n-1})$. For, if $f \neq 0$ is orthogonal to all these elements then

\[
\begin{aligned}\|f-h\|_{2}^{2}&=\|f\|_{2}^{2}-(f,h)-(h,f)+\|h\|_{2}^{2}\\&=\|f\|_{2}^{2}+\|h\|_{2}^{2}\geq\|f\|_{2}^{2}>0\end{aligned}
\]

whenever $h$ is a finite linear combination of members of $\bigcup_{k=0}^{\infty}\mathcal{H}_k$; but, by Corollary \ref{cor:c4-2-3} part (ii), this is impossible. Thus, if $f\in L^2(\Sigma_{n-1})$ then there exists a unique representation,

\begin{equation}\label{eq:c4-2-5}
f=\sum_{k=0}^{\infty}Y^{(k)},
\end{equation}

where the series on the right converges to $f$ in the $L^2$ norm and $Y^{(k)} \in \mathcal{H}_k$; in fact, $Y^{(k)} = b_1^{(k)} Y_1^{(k)} + b_2^{(k)} Y_2^{(k)} + \cdots + b_{a_k}^{(k)} Y_{a_k}^{(k)}$ with $b_j^{(k)} = (f, Y_j^{(k)})$, $j = 1, 2, \ldots, a_k$. When $n = 2$ the series in \eqref{eq:c4-2-5} is the Fourier series of $f$. In this case

\[
Y_{1}^{(k)}(e^{i\theta})=\frac{1}{\sqrt{\pi}}\cos k\theta
\]

and

\[
Y_{2}^{(k)}(e^{i\theta})=\frac{1}{\sqrt{\pi}}\sin k\theta
\]

form an orthonormal basis for $\mathscr{H}_{k}$.

The functions $Y_1^{(k)}$ play a special role in the theory of Fourier series. For example, the Abel means $u(r, \theta) = \sum_{k=-\infty}^{\infty} c_k r^{|k|} e^{ik\theta}$, $0 \leq r < 1$, of the Fourier series of an integrable and periodic function $f$ have simple expressions in terms of the $Y_1^{(k)\prime}$'s and the function $f$. More precisely, using the basic trigonometric identity

\begin{equation}\label{eq:c4-2-6}
\cos\left(\theta-\phi\right)=\cos\phi\cos\theta+\sin\phi\sin\theta
\end{equation}

we find that

\[
u(r,\theta)=\int\limits_{0}^{2\pi}f(\phi)p(r,\theta-\phi)d\phi,
\]

where

\begin{equation}\label{eq:c4-2-7}
p(r,\theta)=\frac{1}{\pi}\left\{\frac{1}{2}+\sum_{k=1}^{\infty}r^{k}\cos k\theta\right\}=\frac{1}{2\pi}\frac{1-r^{2}}{1-2r\cos\theta+\overline{r^{2}}}
\end{equation}

is the Poisson kernel for the unit circle (see Theorem \ref{thm:c2-1-9} in Chapter II). We shall now show that we have a similar situation in higher dimensions. In order to avoid the somewhat exceptional cases of n = 1 and n = 2 (which were anyway treated above), we shall assume that n > 2 throughout the rest of this section.

Let us fix a point $x'$ in $\Sigma_{n-1}$ and consider the linear functional $L$ on $\mathcal{H}_k$ that assigns to each $Y$ in $\mathcal{H}_k$ the value $Y(x')$. By the self-duality of the finite dimensional inner product space $\mathcal{H}_k$ there exists a unique spherical harmonic $Z_x^{(k)}$ such that

\[
L(Y)=Y(x^{\prime})=\int\limits_{\Sigma_{n-1}}Y(t^{\prime})Z_{x^{\prime}}^{(k)}(t^{\prime})dt^{\prime}
\]

for all $Y \in \mathcal{H}_k$. This function $Z_x^{(k)}$ is called the zonal harmonic\index{zonal harmonic@zonal harmonic} of degree k with pole $x'$. We now establish some elementary but basic properties of zonal harmonics.

\begin{lemma}
\label{lem:c4-2-8}
(a) If $\{Y_1, Y_2, \ldots, Y_{a_k}\}$ is an orthonormal basis of $\mathcal{H}_k$ then $Z_{x'}^{(k)}(t') = \sum_{m=1}^{a_k} \overline{Y_m(x')} Y_m(t')$;
(b) $Z_{x'}$ is real-valued and $Z_{x'}^{(k)}(t') = Z_{t'}^{(k)}(x')$;
(c) If $\rho$ is a rotation then $Z_{\rho x'}^{(k)}(\rho t') = Z_{x'}^{(k)}(t')$.
\end{lemma}

\textsc{Proof.}\ Since $\{Y_1, Y_2, \ldots, Y_{a_k}\}$ is an orthonormal basis of $\mathcal{H}_k$ we must have

\[
Z_{x^{\prime}}^{(k)}\;=\;\sum_{m=1}^{a_{k}}\left(Z_{x^{\prime}}^{(k)},Y_{m}\right)Y_{m}.
\]

But, by the defining property of zonal harmonics.

\[
\left(Z_{x^{\prime}}^{(k)},Y_{m}\right)=\int\limits_{\Sigma_{n-1}}\overline{{Y_{m}(t^{^{\prime}})}}Z_{x^{\prime}}^{(k)}(t^{\prime})d t^{\prime}=\overline{{Y_{m}(x^{\prime})}}
\]

and part (a) is proved. Our derivation of the formula for the dimension of $\mathcal{H}_{k}$ was independent of whether we considered the space $\mathcal{H}_{k}$ to consist of real- or complex-valued functions. Consequently, we can choose the orthonormal basis in (a) to consist of real valued functions. If we do this, we see immediately that $Z_{x}^{(k)}$ must be real-valued. Part (b) is now an immediate consequence of part (a).

Letting $w' = \rho t'$ we have, for $Y \in \mathcal{H}_k$,

\[
\begin{aligned}\int\limits_{\Sigma_{n-1}}Z^{(k)}_{\rho x^{\prime}}(\rho t^{\prime})Y(t^{\prime})dt^{\prime}&=\int\limits_{\Sigma_{n-1}}Z^{(k)}_{\rho x^{\prime}}(w^{\prime})Y(\rho^{-1}w^{\prime})dw^{\prime}\\&=Y(\rho^{-1}(\rho x^{\prime}))=Y(x^{\prime}).\end{aligned}
\]

Thus, by the uniqueness of the representation of linear functionals we must have $Z_{\rho x}^{(k)}(\rho t') = Z_{x'}^{(k)}(t')$ and part (c) is proved.

\begin{corollary}
\label{cor:c4-2-9}
(a) $Z_{x'}^{(k)}(x') = a_{k} \omega_{n-1}^{-1}$ for all $x' \in \Sigma_{n-1}$, where $a_k$ is the dimension of $\mathcal{H}_k$ and $\omega_{n-1}$ is the surface area of $\Sigma_{n-1}$;
(b) $\sum_{m=1}^{a_k} |Y_m(x')|^2 = a_k \omega_{n-1}^{-1}$ for all $x' \in \Sigma_{n-1}$, independently of the orthonormal basis $\{Y_1, Y_2, \ldots, Y_{a_k}\}$ of $\mathcal{H}_k$;
(c) $|Z_{t'}^{(k)}(x')| \leq a_{k} \omega_{n-1}^{-1}$ for all $x'$ and $t'$ in $\Sigma_{n-1}$.
\end{corollary}

\textsc{Proof.}\ Suppose $x'_1$ and $x'_2$ are two points of $\Sigma_{n-1}$. We can then find a rotation $\rho$ such that $\rho x'_1 = x'_2$. By part (c) of Lemma \ref{lem:c4-2-8}, therefore, we must have

\[
Z_{x_{2}^{\prime}}^{(k)}(x_{2}^{\prime})=Z_{x_{1}^{\prime}}^{(k)}(x_{1}^{\prime}).
\]

Consequently, $Z_x^{(k)}(x')$ is a constant, independent of $x' \in \Sigma_{n-1}$. From part (a) of Lemma \ref{lem:c4-2-8} we see that this constant, $c$, must equal $\sum_{m=1}^{a_k} |Y_m(x')|^2$, where $\{Y_1, Y_2, \ldots, Y_{a_k}\}$ is an orthonormal basis of $\mathcal{H}_k$. But then

\[
\begin{aligned}a_{k}&=\sum_{m=1}^{a_{k}}\int\limits_{\Sigma_{n-1}}|Y_{m}(x^{\prime})|^{2}dx^{\prime}=\int\limits_{\Sigma_{n-1}}\sum_{m=1}^{a_{k}}|Y_{m}(x^{\prime})|^{2}dx^{\prime}\\&=\int\limits_{\Sigma_{n-1}}c dx^{\prime}=c\omega_{n-1}.\end{aligned}
\]

Thus, $c = a_{k} \omega_{n-1}^{-1}$ which proves both parts (a) and (b).

In order to show part (c) we first observe that it follows from the defining property of zonal harmonics that for $t'$ and $x'$ in $\Sigma_{n-1}$

\begin{equation}
\tag{i}
Z_{t^{^{\prime}}}^{(k)}(x^{\prime})=\intop_{\Sigma_{n-1}}Z_{t^{^{\prime}}}^{(k)}(w^{\prime})Z_{x^{^{\prime}}}^{(k)}(w^{\prime})d w^{\prime}.
\end{equation}

On the other hand, if $\{Y_1, Y_2, \ldots, Y_{a_k}\}$ is an orthonormal basis of $\mathcal{H}_k$ it follows from Lemma \ref{lem:c4-2-8}, part (a), and the result we have just established that for all $u' \in \Sigma_{n-1}$

\[
\Vert Z_{u^{\prime}}^{(k)}\Vert_{2}^{2}=\int_{\Sigma_{n-1}}\vert Z_{u^{\prime}}^{(k)}(w^{\prime})\vert^{2}d w^{\prime}=\sum_{m=1}^{a_{k}}\vert Y_{m}(u^{\prime})\vert^{2}=a_{k}\omega_{n-1}^{-1}.
\]

Thus, (i), Schwarz's inequality and this last equality give us

\[
\left|Z_{t^{\prime}}^{(k)}(x^{\prime})\right|\leq\|Z_{t^{\prime}}^{(k)}\|_{2}\|Z_{x^{\prime}}^{(k)}\|_{2}=\sqrt{a_{k}\omega_{n-1}^{-1}}\sqrt{a_{k}\omega_{n-1}^{-1}}=a_{k}\omega_{n-1}^{-1}
\]

and the last part of the corollary is established. \footnote{The reader should observe that when $n = 2$ Lemma \ref{lem:c4-2-8} and Corollary \ref{cor:c4-2-9} reduce to familiar elementary properties of the trigonometric functions. For example, Lemma \ref{lem:c4-2-8} (a) is then essentially the identity \eqref{eq:c4-2-6} and Corollary \ref{cor:c4-2-9} (b) extends the equality $\cos^2 \theta + \sin^2 \theta = 1$.}

From formula \eqref{eq:c4-2-7} we see how the Poisson kernel for the unit circle can be expressed in terms of zonal harmonics. We now show that such an expression is valid in $n$-dimensions. We recall (see Theorem \ref{thm:c2-1-9} of Chapter II) that the Poisson kernel for the unit sphere in $E_n$ is given by the formula

\[
p(t^{\prime},x)=\frac{1}{\omega_{n-1}}\frac{1-|x|^{2}}{|x-t^{\prime}|^{n}}
\]

for $0 \leq |x| < 1 = |t'|$.

\begin{theorem}
\label{thm:c4-2-10}
Let $x = rx'$ with $r = |x| < 1$ then
\[
p(t^{\prime},x)=\sum_{k=0}^{\infty}r^{k}Z_{x^{\prime}}^{(k)}(t^{\prime})=\sum_{k=0}^{\infty}r^{k}Z_{t^{\prime}}^{(k)}(x^{\prime})
\]

for all $t' \in \Sigma_{n-1}$.
\end{theorem}

\textsc{Proof.}\ It follows from Corollary \ref{cor:c4-2-9} (c) and the fact that

\[
a_{k}=d_{k}-d_{k-2}=\{(n+2k-2)/k\}\binom{n+k-3}{k-1}
\]

is less than or equal to a constant times $k^{(n-2)}$ that the series

\[
q(t^{\prime},x)=\sum_{k=0}^{\infty}r^{k}Z_{t^{\prime}}^{(k)}(x^{\prime})
\]

converges uniformly in each closed region $\{x \in E_n; |x| \leq r_0 < 1\}$. Now suppose $u(t') = \sum_{j=0}^m Y_j(t')$ is a finite linear combination of spherical harmonics (where $Y_j \in \mathcal{H}_j$). Then by Theorem \ref{thm:c2-1-10} of Chapter II, the definition of spherical harmonics, and the uniqueness of the solution of the Dirichlet problem

\[
\sum_{j=0}^{m}|x|^{j}Y_{j}(x^{\prime})=u(x)=\int\limits_{\Sigma_{n-1}}u(t^{\prime})p(t^{\prime},x)dt^{\prime}
\]

must be the continuous function on $|x| \leq 1$ that is harmonic for $|x| < 1$ and equal to $u(x')$ for $|x'| = 1$. But by Lemma \ref{lem:c4-2-8} (b), the defining property of zonal harmonics, and the orthogonality relation Corollary \ref{cor:c4-2-4},

\[
\begin{aligned}\int\limits_{\Sigma_{n-1}}u(t^{\prime})q(t^{\prime},x)dt^{\prime}&=\sum_{j=0}^{m}\int\limits_{\Sigma_{n-1}}Y_{j}(t^{\prime})q(t^{\prime},x)dt^{\prime}\\&=\sum_{j=0}^{m}\left\{\sum_{k=0}^{\infty}|x|^{k}\int\limits_{\Sigma_{n-1}}Y_{j}(t^{\prime})Z_{t^{\prime}}^{(k)}(x^{\prime})dt^{\prime}\right\}\\&=\sum_{j=1}^{m}|x|^{j}Y_{j}(x^{\prime})=u(x).\end{aligned}
\]

Thus, $\int_{\Sigma_{n-1}} [p(t', x) - q(t', x)] u(t') \, dt' = 0$ for all finite linear combinations of spherical harmonics. Since the latter are dense in $L^2(\Sigma_{n-1})$ (Corollary \ref{cor:c4-2-3}, part (ii)) and both $p(t', x)$ and $q(t', x)$ are continuous functions of $t'$ it follows that $p(t', x) = q(t', x)$ and our theorem is established.

The zonal harmonics can be characterized by a simple geometrical property. In order to describe this property we define a parallel of $\Sigma_{n-1}$ orthogonal to a point $e$ in $\Sigma_{n-1}$ to be the intersection of the unit sphere with a hyperplane that is perpendicular to the line determined by the origin and $e$. If $\rho$ is a rotation leaving $e$ fixed it then follows from Lemma \ref{lem:c4-2-8} (c) that $Z_e^{(k)}(x') = Z_e^{(k)}(\rho x')$ for all $x'$ in $\Sigma_{n-1}$. But such a rotation $\rho$ must clearly map a parallel orthogonal to $e$ into itself; moreover, given two points, $x_1'$ and $x_2'$, on this parallel it follows from the two-dimensional case that there exists a rotation $\rho$ that leaves $e$ fixed and such that $\rho x_1' = x_2'$ (for $\rho$ can be chosen so that it leaves the orthogonal complement of the plane spanned by $x_1'$ and $x_2'$ fixed). Consequently, the zonal harmonic $Z_e^{(k)}$ with pole $e$ is constant on the parallels of $\Sigma_{n-1}$ orthogonal to $e$. We shall show that this property characterizes the zonal harmonics up to constant multiples. We observe that such functions are invariant under the rotations which leave $e$ fixed. To utilize this remark we also need the following result.

\begin{lemma}
\label{lem:c4-2-11}
Suppose $P$ is a polynomial on $E_n, n \geq 2$, such that $P(\rho x) = P(x)$ for all rotations $\rho$ and points $x \in E_n$, then there exist constants $c_0, c_1, c_2, \ldots, c_m$ such that
\[
P(x)=\sum_{k=0}^{m}c_{k}(x_{1}^{2}+x_{2}^{2}+\cdots+x_{n}^{2})^{k}.
\]

\end{lemma}

\textsc{Proof.}\ We can always write $P(x) = \sum_{l=0}^{j} P_l(x)$, where $P_l$ is homogeneous of degree $l$. Then, for each $\varepsilon > 0$ and each rotation $\rho$,

\[
\sum_{l=0}^{j}\varepsilon^{l}P_{i}(x)=P(\varepsilon x)=P(\varepsilon\rho x)=\sum_{l=0}^{j}\varepsilon^{l}P_{l}(\rho x).
\]

Consequently, we must have $P_i(\rho x) = P_i(x)$ for $l = 0, 1, \ldots, j$. If we let $F(x) = |x|^{-1}P_i(x)$ then $F$ is a homogeneous function of degree 0 that is invariant under the rotation group (that is, $F(\rho x) = F(x)$ for all rotations $\rho$). But this implies that $F$ is a constant function: $F(x) = c'_i$ for all $x$ in $E_n$. Thus, $P_i(x) = c'_i|x|^l$. Since $P_i$ is a polynomial, however, $l$ must be even when $c'_i \neq 0$. Thus, $P(x) = \sum_{k=0}^{m} c_k|x|^{2k}$, where $c_k = c'_{2k}$ for $k = 0, 1, 2, \ldots, m$ and $m$ is the largest integer less than or equal to $j/2$.

\begin{theorem}
\label{thm:c4-2-12}
Suppose $e$ is a point of $\Sigma_{n-1}$, then $Y \in \mathcal{H}_k$ is constant on parallels of $\Sigma_{n-1}$ orthogonal to $e$ if and only if there exists a constant $c$ such that $Y = c Z_e^{(k})$.
\end{theorem}

\textsc{Proof.}\ We have already shown that zonal harmonics have this property. Suppose, therefore, that $Y$ is constant on parallels of $\Sigma_{n-1}$ orthogonal to $e$. If $e_1 = (1, 0, \ldots, 0) \in \Sigma_{n-1}$ and $\tau$ is a rotation satisfying $e = \tau e_1$, then the spherical harmonic $W$ whose values are $W(x') = Y(\tau x')$ is constant on parallels of $\Sigma_{n-1}$ orthogonal to $e_1$. If we show that $W = c Z_{e_1}^{(k)}$ then, by Lemma \ref{lem:c4-2-8} (c),

\[
\begin{array}{r}{Y(y^{\prime})=W(\tau^{-1}y^{\prime})=c Z_{e_{1}}^{(k)}(\tau^{-1}y^{\prime})=c Z_{\tau^{-1}e}^{(k)}(\tau^{-1}y^{\prime})=c Z_{e}^{(k)}(y^{\prime})}\end{array}
\]

for all $y' \in \Sigma_{n-1}$. Consequently, it suffices to show that $W = cZ_{e_1}^{(k)}$.

Let $P(x) = |x|^k W(x/|x|)$ for $x \neq 0$ and $P(0) = 0$. Then, if $\rho$ is a rotation leaving $e_1$ fixed then $P(\rho x) = P(x)$ for all $x \in E_n$ and, also, all polynomials of the form $x_1^m$, $m = 0, 1, 2, \ldots$, are invariant under the action of $\rho$. Consequently, if we write

\[
P(x)=\sum_{j=0}^{k}x_{1}^{k-j}P_{j}(x_{2},x_{3},\ldots,x_{n})
\]

it follows that $P_0, P_1, \ldots, P_k$ are invariant under the action of $\rho$ (observe that $\rho(x_1, x_2, \ldots, x_n) = (x_1, x_2', \ldots, x_n')$ and the mapping $(x_2, \ldots, x_n) \to (x_2', \ldots, x_n')$ is a rotation of the $(n-1)$-dimensional space $x_1 = 0$; moreover, by varying $\rho$ over all rotations leaving $e_1$ fixed we obtain all rotations of this $(n-1)$-dimensional space in this way). By Lemma \ref{lem:c4-2-11}, therefore, $P_j$ is zero if $j$ is odd while $P_j(x_2, \ldots, x_n) = c_j(x_2^2 + \cdots + x_n^2)^{j/2}$ if $j$ is even. Letting $R = (x_2^2 + \cdots + x_n^2)^{1/2}$, therefore, we have shown that

\begin{equation}
\tag{i}
\begin{array}{r}{P(x)=c_{0}x_{1}^{k}+c_{2}x_{1}^{k-2}R^{2}+\cdots+c_{2l}x_{1}^{k-2l}R^{2l}.}\end{array}
\end{equation}

Because W is a surface spherical harmonic the polynomial P must be a solid spherical harmonic. Thus,

\[
0=\Delta P(x)=\sum_{j=0}^{l-1}[c_{2j}\alpha_{j}+c_{2(j+1)}\beta_{j}]x_{1}^{k-2(j+1)}R^{2j},
\]

where $\alpha_j = (k - 2j)(k - 2j - 1)$ and $\beta_j = 2(j + 1)(n + 2j - 1)$. Hence, $c_{2(j+1)} = -(\alpha_j/\beta_j)c_{2j}$ for $j = 0, 1, 2, \ldots, l - 1$. But this asserts that all the coefficients $c_0, c_2, \ldots, c_{2l}$ are determined by $c_0$; consequently, any two nonzero harmonic polynomials having form (i) must be constant multiples of each other. On the other hand, we have shown that any homogeneous polynomial of degree $k$ whose restriction to $\Sigma_{n-1}$ is constant on parallels of $\Sigma_{n-1}$ orthogonal to $e_1$ must have form (i). Since, $Z_{e_1}^{(k)}(x/|x|)|x|^{k}$ has this property and, moreover, is harmonic it follows that $W(x') = P(x') = cZ_{e_1}^{(k)}(x')$ for all $x' \in \Sigma_{n-1}$, and the theorem is proved.

\begin{corollary}
\label{cor:c4-2-13}
Suppose $F_{y}(x')$ is defined for all $x', y' \in \Sigma_{n-1}$ and
(a) $F_y$. is a spherical harmonic of degree k for each $y' \in \Sigma_{n-1}$,
(b) if $\rho$ is a rotation then $F_{\rho y^{\prime}}(\rho x^{\prime}) = F_{y^{\prime}}(x^{\prime})$, then there exists a constant $c$ such that $F_{y^{\prime}}(x^{\prime}) = c Z_{y^{\prime}}^{(k)}(x^{\prime})$ for all $x^{\prime}$, $y^{\prime}$ in $\Sigma_{n-1}$.
\end{corollary}

\textsc{Proof.}\ Fix $y'$ in $\Sigma_{n-1}$ and let $\rho$ be a rotation leaving $y'$ fixed. Then, using (b) we have

\[
F_{y^{\prime}}(x^{\prime})=F_{\rho y^{\prime}}(\rho x^{\prime})=F_{y^{\prime}}(\rho x^{\prime})
\]

for all $x'$ in $\Sigma_{n-1}$. But this means, because of assumption (a), that $F_{y'}$ is a spherical harmonic that is constant on parallels of $\sum_{n-1}$ orthogonal to $y'$. By theorem 2.12, there exists $c(y')$ such that $F_{y'} = c(y')Z_{y}^{(k)}$. The corollary will be proved, therefore, if we can show that $c(y_1') = c(y_2')$ whenever $y_1'$ and $y_2'$ belong to $\Sigma_{n-1}$. In order to do this we consider a rotation $\sigma$ such that $\sigma y_1' = y_2'$. Then, using assumption (b), $c(y_2)Z_{y_2}^{(k)}(\sigma x') = F_{y_2}(\sigma x') = F_{\sigma y_1}(\sigma x') = F_{y_1}(x') = c(y_1')Z_{y_1}^{(k)}(x')$. On the other hand, by Lemma \ref{lem:c4-2-8} (c),

\[
Z_{y_{1}^{\prime}}^{(k)}(x^{\prime})=Z_{\sigma y_{1}^{\prime}}^{(k)}(\sigma x^{\prime})=Z_{y_{2}^{\prime}}^{(k)}(\sigma x^{\prime}).
\]

Consequently, $c(y_{1}^{\prime}) = c(y_{2}^{\prime})$.

The zonal harmonics have a particularly simple expression in terms of the ultraspherical\index{ultraspherical (Gegenbauer) polynomials@ultraspherical (Gegenbauer) polynomials} (or Gegenbauer) polynomials $P_{\kappa}^{\lambda}$. The latter can be defined in terms of a generating function. If we write

\[
(1-2rt+r^{2})^{-\lambda}=\sum_{k=0}^{\infty}P_{k}^{\lambda}(t)r^{k},
\]

where $0 \leq |r| < 1$, $|t| \leq 1$ and $\lambda > 0$, then the coefficient $P_k^\lambda(t)$ is called the ultraspherical polynomial of degree $k$ associated with $\lambda$. We shall derive some of the most elementary properties of the functions $P_k^\lambda$. Putting $r = 0$ we see that

(i)

\[
P_{0}^{\lambda}(t)\equiv1.
\]

Since

\[
\begin{aligned}2r\lambda\sum_{k=0}^{\infty}P_{k}^{\lambda+1}(t)r^{k}&=2r\lambda(1-2rt+r^{2})^{-\lambda-1}\\&=\frac{d}{dt}(1-2rt+r^{2})^{-\lambda}\\&=\sum_{k=0}^{\infty}r^{k}\frac{d}{dt}P_{k}^{\lambda}(t)\end{aligned}
\]

we have

\begin{equation}
\tag{ii}
\frac{d}{d t}P_{k}^{\lambda}(t)\:=\:2\lambda P_{k-1}^{\lambda+1}(t)\quad\text{for }k\geq1.
\end{equation}

From (i) and (ii) we have $(d/dt)P_1^\lambda(t) = 2\lambda$; thus, $P_1^\lambda(t)$ is a polynomial of degree 1. An induction argument using property (ii) gives us

(iii) $P_{k}^\lambda(t)$ is a polynomial in $t$ of (precisely) degree $k$. It follows from this last property that the polynomials $1, t, t^2, \ldots, t^k, \ldots$ are obtainable as finite linear combinations of $P_0, P_1, \ldots, P_k, \ldots$ Thus, by the Weierstrass approximation theorem,

(iv) The finite linear combinations of the polynomials $P_k^\lambda(t), k = 0, 1, \ldots$, are uniformly dense in the space of continuous functions on the closed interval [-1, 1].

Since

\[
\sum_{k=0}^{\infty}P_{k}^{\lambda}(-t)r^{k}=(1+2rt+r^{2})^{-\lambda}=\sum_{k=0}^{\infty}P_{k}^{\lambda}(t)(-r)^{k}
\]

we must have

\begin{equation}
\tag{v}
P_{k}^{\lambda}(-t)=(-1)^{k}P_{k}^{\lambda}(t)\quad\text{for }k\geq0.
\end{equation}

The following representation of the zonal harmonics is an easy consequence of these properties:

\begin{theorem}
\label{thm:c4-2-14}
If $n > 2$ is an integer, $\lambda = (n - 2)/2$ and $k = 0, 1, 2, \ldots$ then there exists a constant $c_{k,n}$ such that
\[
\begin{array}{r}{Z_{y^{\prime}}^{(k)}(x^{\prime})=c_{k,n}P_{k}^{\lambda}(x^{\prime}{\cdot}y^{\prime})}\end{array}
\]

for all $x', y' \in \Sigma_{n-1}$.
\end{theorem}

\textsc{Proof.}\ For $y'$ in $\Sigma_{n-1}$ we let $F_y(x) = |x|^{k}P_k^\lambda(x \cdot y'/|x|)$ for $x$ in $E_n$. If we can show that $F_y'(x')$ satisfies the hypotheses of Corollary \ref{cor:c4-2-13} the theorem would then be immediate. By properties (iii) and (v), if $k$ is even, say $k = 2m$, then $P_k^\lambda(t)$ must have the form $P_k^\lambda(t) = \sum_{j=0}^{m} d_{2j} t^{2j}$ while, if $k$ is odd, say $k = 2m + 1$, $P_k^\lambda(t) = \sum_{j=0}^{m} d_{2j+1} t^{2j+1}$. Thus, in the former case,

\[
|x|^{k}P_{k}^{\lambda}\biggl(\frac{x\cdot y^{\prime}}{|x|}\biggr)\;=\;\sum_{j=0}^{m}d_{2j}(x_{1}^{2}+\cdots+x_{n}^{2})^{m-j}(x_{1}y_{1}^{\prime}+\cdots+x_{n}y_{n}^{\prime})^{2j},
\]

and, in the second case

\[
|x|^{k}P_{k}^{\lambda}\biggl(\frac{x\cdot y^{\prime}}{|x|}\biggr)=\sum_{j=0}^{m}d_{2j+1}(x_{1}^{2}+\cdots+x_{n}^{2})^{m-j}(x_{1}y_{1}^{\prime}+\cdots+x_{n}y_{n}^{\prime})^{2j+1}.
\]

In either case, $F_{y^{\prime}}(x)$ is a (nonzero) homogeneous polynomial of degree k. If $\rho$ is a rotation then, since it preserves the inner product,

\[
F_{\rho y^{\prime}}(\rho x^{\prime})=P_{k}^{\lambda}(\rho x^{\prime}\cdot\rho y^{\prime})=P_{k}^{\lambda}(x^{\prime}\cdot y^{\prime})=F_{y^{\prime}}(x^{\prime})
\]

for all points $x'$, $y'$ of $\Sigma_{n-1}$. Thus, property (b) of Corollary \ref{cor:c4-2-13} is satisfied; consequently, we only need to show that $F_{y'}$ is harmonic. We have

already observed (see Examples 3 and 5 at the beginning of Section 1 in Chapter II) that $|x - x_0|^{2-n}$ is a harmonic function of x in the region $E_n - \{x_0\}$. In particular,

\begin{equation}\label{eq:c4-2-15}
\begin{aligned}s^{2-n}\Bigg|x-\frac{y^{\prime}}{s}\Bigg|^{2-n}&=\bigg[1-2s|x|\bigg(\frac{x}{|x|}\cdot y^{\prime}\bigg)\;+\;(s|x|)^{2}\bigg]^{-\lambda}\\&=\sum_{k=0}^{\infty}s^{k}|x|^{k} P_{k}^{\lambda}\bigg(\frac{x}{|x|}\cdot y^{\prime}\bigg)\end{aligned}
\end{equation}

is harmonic in the region $\mathcal{R} = \mathcal{R}_s = \{x \in E_n; 0 < |x| < 1/s\}$, for $s \neq 0$ and $y'$ in $\Sigma_{n-1}$ fixed. But this implies that each of the coefficients $F_{y'}(x) = |x|^k P_k^k(x, y'/|x|))$ must be a harmonic function of $x$. This can be seen, for example, by taking the mean value of both the left and right sides of equality \eqref{eq:c4-2-15} on spheres about $x$ that lie within the region $\mathcal{R}_s$. Since the left side satisfies the mean-value property (Theorem \ref{thm:c2-1-1} of Chapter II) for $0 < s \leq s_0 < \infty$, the same must be true for each of the coefficients $F_{y'}(x) = |x|^k P_k^k(x, y'/|x|)$. From this and Theorem \ref{thm:c2-1-7} of Chapter II we can now deduce the desired result that $F_{y'}$ is harmonic.

The ultraspherical polynomials can be used to obtain some important orthogonal expansions of functions. The basic orthogonality property needed for carrying out such expansions is an easy consequence of the last theorem.

\begin{corollary}
\label{cor:c4-2-16}
The polynomials $P_{k}^{(n-2)/2}(t)$, $k = 0, 1, 2, \ldots$, are mutually orthogonal with respect to the inner product
\[
(f,g)=\int\limits_{-1}^{1}f(t)g(t)(1-t^{2})^{(n-3)/2}d t.
\]

\end{corollary}

\textsc{Proof.}\ Let $e = (1, 0, \ldots, 0)$ and, for $x' \in \Sigma_{n-1}$, define $\theta, 0 \leq \theta \leq \pi$, by letting $e \cdot x' = \cos \theta$. We can evaluate integrals over $\Sigma_{n-1}$ by first integrating over the parallel $L_\theta = \{x' \in \Sigma_{n-1}; e \cdot x' = \cos \theta\}$ orthogonal to $e$, thus obtaining a function of $\theta$, $0 \leq \theta \leq \pi$, which can then be integrated over the interval $[0, \pi]$. The measure of $L_\theta$ is $\omega_{n-2}(\sin \theta)^{n-2}$ (the surface area of a sphere in $E_{n-2}$ of radius $\sin \theta$). Using this fact together with Corollary \ref{cor:c4-2-4}, Theorems \ref{thm:c4-2-12} and 2.14 we thus have

\[
\begin{aligned}0&=\int\limits_{\Sigma_{n-1}}Z_{e}^{l}(x^{\prime})Z_{e}^{k}(x^{\prime})dx^{\prime}\\&=c_{l,n}c_{k,n}\omega_{n-2}\int\limits_{0}^{\pi}P_{l}^{(n-2)/2}\left(\cos\theta\right)P_{k}^{(n-2)/2}\left(\cos\theta\right)(\sin\theta)^{n-2}d\theta\\&=\omega_{n-2}c_{l,n}c_{k,n}\int\limits_{-1}^{1}P_{l}^{(n-2)/2}(t)P_{k}^{(n-2)/2}(t)(1-t^{2})^{(n-3)/2}dt.\end{aligned}
\]

From this result and property (iv) of the ultraspherical polynomials we obtain

\begin{corollary}
\label{cor:c4-2-17}
The polynomials $P_k^{(n-2)/2}, k = 0, 1, 2, \ldots$, form an orthogonal basis for the space $L^2([-1,1];(1-t^2)^{(n-3)/2} dt).\footnote{Corollaries \ref{cor:c4-2-16} and \ref{cor:c4-2-17} are valid for all $\lambda > 0$ and not just $\lambda = (n - 2)/2$. Different arguments, however, are needed to establish them.}$
Having completed our study of the spherical harmonics, mainly from the point of view of functions on the sphere, we shall now show how they can be used in the Fourier analysis of functions on $E_n$. More precisely, we shall study the subspaces $\mathfrak{G}_k$ of $L^2(E_n)$. We recall that $\mathfrak{G}_k$ was defined to be the space of all linear combinations of functions of the form $f(r)P(x)$, where $f$ ranges over the radial functions and $P$ over the solid spherical harmonics of degree $k$, in such a way that $f(r)P(x)$ belongs to $L^2(E_n)$.
\end{corollary}

We now state and prove the following lemma as preparatory to the development in Section 3.

\begin{lemma}
\label{lem:c4-2-18}
The direct sum decomposition $L^2(E_n) = \sum_{k=0}^{\infty} \oplus \mathfrak{G}_k$ holds in the sense that
(a) each subspace $\mathfrak{Q}_{k}$ is closed;
(b) the $Q_{k}$ are mutually orthogonal;
(c) every element of $L^2(E_n)$ is a limit of finite linear combinations of elements belonging to the spaces $\mathfrak{H}_k$. In addition, the Fourier transform maps $\mathfrak{H}_k$ into itself.
\end{lemma}

\textsc{Proof.}\ The space of solid spherical harmonics of degree k, being isomorphic to $\mathcal{H}_k$, has (finite) dimension $a_k = d_k - d_{k-2}$. Let $P_1, P_2, \ldots, P_{a_k}$ be an orthonormal basis of this space (where the inner product is that which is inherited from $L^2(\Sigma_{n-1}))$. Then, clearly, every element of $\mathfrak{S}_k$ can be written in the form $\sum_{j=1}^{a_k} f_j(r) P_j(x)$. Moreover,

\[
\int\limits_{E_{n}}|f(x)|^{2}\,dx=\sum_{j=1}^{a_{k}}\int\limits_{0}^{\infty}|f_{j}(r)|^{2}r^{2k+n-1}\,dr.
\]

From this conclusion (a) is obvious. The mutual orthogonality of the $\mathfrak{S}_{\kappa}$'s follows immediately from the orthogonality of the spherical harmonics (Corollary \ref{cor:c4-2-4}) and by an integration involving polar coordinates. To prove the completeness property (c), it suffices to show that, if a function of $L^{2}(E_{n})$ is orthogonal to all the $\mathfrak{S}_{\kappa}$'s then it vanishes almost everywhere. By the completeness of the spherical harmonics on the sphere, however, such a function must vanish a.e. on almost every sphere centered at the origin, which gives the desired conclusion.

The action of the Fourier transform on $\mathfrak{H}_{k}$ will be the subject matter of the following section. An elementary argument, however, enables us to

prove at this point that the Fourier transform preserves $\mathfrak{S}_k$. It suffices, for this purpose, to consider $f\in L^1(E_n)\cap L^2(E_n)$ having the form $f(u)=f_0(\rho)P(u)=\rho^k f_0(\rho)Y(u')$ with $Y\in\mathcal{H}_k$, where $\rho=|u|$ and $u=\rho u'$. Since the finite linear combinations of such functions are dense in $\mathfrak{S}_k$, this last space will be invariant under the action of the Fourier transform provided $\hat{f}\in\mathfrak{S}_k$ whenever $f$ has the above form. Letting $r=|x|$ and $x=rx'$ we then obtain

\[
\begin{aligned}&\hat{f}(x)=\int_{E_{n}}e^{-2\pi i x\cdot u}f(u)\;d u\\&\qquad=\int_{0}^{\infty}f_{0}(\rho)\rho^{k+n-1}\Biggl\{\int_{\Sigma_{n-1}}e^{-2\pi i r\rho x^{\prime}\cdot u^{\prime}}Y(u^{\prime})\;d u^{\prime}\Biggr\}d\rho.\end{aligned}
\]

If we can show that there exists a function $\varphi$, defined on $[0, \infty)$, such that

\begin{equation}\label{eq:c4-2-19}
\intop_{\Sigma_{n-1}}e^{-2\pi\mathrm{i}s x^{\prime}\cdot u^{\prime}}Y(u^{\prime})\;d u^{\prime}\;=\;\varphi(s)Y(x^{\prime})
\end{equation}

for $s \geq 0$ then

\[
\hat{f}(x)=\left\{\int\limits_{0}^{\infty}f_{0}(\rho)\varphi(r\rho)\rho^{k+n-1}d\rho\right\}Y(x^{\prime})
\]

and, consequently, $\hat{f} \in \mathfrak{H}_{k}$.

By the defining property of the zonal harmonic $Z_{u'}^{(k)}(v')$ and the property $Z_{u'}^{(k)}(v') = Z_{v'}^{(k)}(u')$ we must have

\[
\begin{aligned}\int\limits_{\Sigma_{n-1}}e^{-2\pi\mathrm{i}s x^{\prime}\cdot u^{\prime}}Y(u^{\prime})\,du^{\prime}&=\int\limits_{\Sigma_{n-1}}e^{-2\pi\mathrm{i}s x^{\prime}\cdot u^{\prime}}\Biggl\{\int\limits_{\Sigma_{n-1}}Y(v^{\prime})Z_{u^{\prime}}^{(k)}(v^{\prime})\,dv^{\prime}\Biggr\}du^{\prime}\\&=\int\limits_{\Sigma_{n-1}}Y(v^{\prime})\Biggl\{\int\limits_{\Sigma_{n-1}}e^{-2\pi\mathrm{i}s x^{\prime}\cdot u^{\prime}}Z_{v^{\prime}}^{(k)}(u^{\prime})\,du^{\prime}\Biggr\}dv^{\prime}.\end{aligned}
\]

If we denote the expression in the last brackets by $F_{x}(v')$ an immediate application of Fubini's theorem and Corollary \ref{cor:c4-2-4} gives us the fact that $F_{x'}$, as a function of $v' \in \Sigma_{n-1}$, is orthogonal to all the spaces $\mathcal{H}_{j}, j \neq k$. But, by \eqref{eq:c4-2-5} this implies that $F_{x'} \in \mathcal{H}_{k}$. From Lemma \ref{lem:c4-2-8}, part (c) and the change of variables $u' = \sigma w',$ we obtain the property $F_{\sigma x'}(\sigma v') = F_{x}(v')$ for all rotations $\sigma$. Thus, by Corollary \ref{cor:c4-2-13} there exists a number $c( = \varphi(s))$ such that $F_{x}(v') = cZ_{x}^{(k)}(v')$ for all $x', v'$ in $\Sigma_{n-1}$. Consequently,

\[
\begin{aligned}\int\limits_{\Sigma_{n-1}}\!e^{-2\pi i s x^{\prime}\cdot u^{\prime}}Y(u^{\prime}) du^{\prime}&=\int\limits_{\Sigma_{n-1}}Y(v^{\prime})F_{x^{\prime}}(v^{\prime}) dv^{\prime}=\int\limits_{\Sigma_{n-1}}Y(v^{\prime})\varphi(s)Z_{x^{\prime}}^{(k)}(v^{\prime}) dv^{\prime}\\&=\varphi(s)Y(x^{\prime})\end{aligned}
\]

and \eqref{eq:c4-2-19} is established.

\section{The Action of the Fourier Transform on the Spaces  $\mathfrak{Q}_{k}$}

We have shown that the Fourier transform of a radial function is a radial function (see Corollary \ref{cor:c4-1-2}). In the two dimensional case we have obtained an explicit formula relating these two radial functions (the special case $k=0$ of Theorem \ref{thm:c4-1-6}). This formula has a natural extension to $n>2$ dimensions and, as is the case when $n=2$, it involves a Bessel function. We shall show that this Bessel function is $J_{(n-2)/2}$. When $n$ is odd ($n-2)/2$ is not an integer; consequently, in this case $J_{(n-2)/2}$ is not one of the functions introduced in the first section. We shall show, however, that there is a natural way of extending the class of Bessel functions we introduced so as to include the ones we will encounter. In order\index{Bessel function@Bessel function!of arbitrary order@of arbitrary order} to do this we establish the following identity (known as the Poisson representation\index{Poisson representation of Bessel functions@Poisson representation of Bessel functions} of Bessel functions):

\begin{lemma}
\label{lem:c4-3-1}
If k is a nonnegative integer then
\[
J_{k}(t)=\frac{(t/2)^{k}}{\Gamma[(2k+1)/2]\Gamma(\frac{1}{2})}\int_{-1}^{1}e^{i t s}(1-s^{2})^{(2k-1)/2}d s.
\]

\end{lemma}

\textsc{Proof.}\ We define $J_k^*$ by letting

\[
J_{k}^{*}(t)=\frac{(t/2)^{k}}{\Gamma[(2k+1)/2]\Gamma(\frac{1}{2})}\int_{-1}^{1}e^{i t s}(1-s^{2})^{(2k-1)/2}d s.
\]

From the definition of $J_0$ (see Section 1) and the change of variables $s = \sin \theta$ we immediately obtain the fact that $J_0^* = J_0$. The lemma will be established, therefore, if we can show that the recursion relation

\begin{equation}\label{eq:c4-3-2}
\frac{d}{d t}\left(t^{-k}G_{k}(t)\right)=-t^{-k}G_{k+1}(t),t\neq0,\qquad k=0,1,2,\ldots,
\end{equation}

is satisfied by both of the sequences $\{J_{k}\}$ and $\{J_{k}^{*}\}$. But,

\[
\begin{aligned}{\frac{d}{d t}}&{{}(t^{-k}J_{k}(t))=-t^{-k}\biggl\{\frac{k}{t}J_{k}(t)-J_{k}^{\prime}(t)\biggr\}}\\ {}&{{}=-t^{-k}\Biggl\{\frac{k}{2\pi t}\intop_{0}^{2\pi}e^{i t\sin\theta}e^{-i k\theta}d\theta-\frac{1}{2\pi}\intop_{0}^{2\pi}\biggl(\frac{d}{d t}e^{i t\sin\theta}\biggr)e^{-i k\theta}d\theta\Biggr\}}\\ {}&{{}=-\frac{t^{-k}}{2\pi}\intop_{0}^{2\pi}\biggl\{i\frac{d}{d\theta}\left[\frac{e^{i t\sin\theta}e^{-i k\theta}}{t}\right]}\\
 {}&{{}\quad+\cos\theta e^{i t\sin\theta}e^{-i k\theta}-i\sin\theta e^{i t\sin\theta}e^{-i k\theta}\biggr\}d\theta}\\ {}&{{}=-\frac{t^{-k}}{2\pi}\intop_{0}^{2\pi}e^{i t\sin\theta}e^{-i(k+1)\theta}d\theta=\left.-t^{-k}J_{k+1}(t)\right..}\\ \end{aligned}
\]

While, integrating by parts and using the fact that

\[
\left(\frac{2k+1}{2}\right)\Gamma\left(\frac{2k+1}{2}\right)=\Gamma\left(\frac{2k+3}{2}\right),
\]

we have

\[
\begin{aligned}\frac{d}{dt}(t^{-k}J^{*}_{k}(t))&=\frac{2^{-k}i}{\Gamma[(2k+1)/2]\Gamma(\frac{1}{2})}\int\limits_{-1}^{1}e^{its}s(1-s^{2})^{(2k-1)/2}ds\\&=\frac{2^{-k}i}{\Gamma[(2k+1)/2]\Gamma(\frac{1}{2})}\int\limits_{-1}^{1}\frac{2it}{(2k+1)}e^{its}\frac{(1-s^{2})^{(2k+1)/2}}{2}ds\\&=-t^{-k}J^{*}_{k+1}(t).\\ \end{aligned}
\]

The integral occurring in Lemma \ref{lem:c4-3-1} is well-defined for all real $k > -\frac{1}{2}$. We therefore obtain a wider class of functions than the one considered in Section 1 if we define the Bessel function $J_k$, where $k$ is a real number greater than $-\frac{1}{2}$, by letting

\[
J_{k}(t)=\frac{(t/2)^{k}}{\Gamma[(2k+1)/2]\Gamma(\frac{1}{2})} \intop_{-1}^{1}e^{i t s}(1-s^{2})^{(2k-1)/2} d s
\]

for t > 0.

Now suppose $f$ is a radial function in $L^{1}(E_{n})$; thus $f(x) = f_{0}(|x|)$ for $a.e.$ $x$ in $E_{n}$, where $\int_{0}^{\infty} |f_{0}(r)|r^{n-1}dr < \infty$. The Fourier transform $\hat{f}$ is then also radial; consequently, $\hat{f}(x) = F_{0}(|x|)$ for all $x \in E_{n}$. Thus, if $r = |x|$, $x = rx^{\prime}$, $s = |u|$ and $u = su^{\prime}$, we have $F_{0}(r) = \hat{f}(x) = \int_{E_{n}} f(u)e^{-2\pi i x \cdot u} du = \int_{0}^{\infty} f_{0}(s)\{\int_{\Sigma_{n-1}} e^{-2\pi i s(x^{\prime}\cdot u^{\prime})} du^{\prime}\}s^{n-1}ds$. We can evaluate the inner integral by the method used in the proof of Corollary \ref{cor:c4-2-16}: First we integrate over the parallel $L_{\theta} = \{u^{\prime} \in \Sigma_{n-1}; x^{\prime} \cdot u^{\prime} = \cos \theta\}$ orthogonal to $x^{\prime}$ obtaining a function of $\theta$, $0 \leq \theta \leq \pi$, which we then integrate over the interval $[0, \pi]$. Thus, since $e^{-2\pi i s(x^{\prime}\cdot u^{\prime})} = e^{-2\pi i s\cos\theta}$ is constant on $L_{\theta}$ and the measure of this parallel is $\omega_{n-2}(\sin \theta)^{n-2} = (2\pi^{(n-1)/2}/\Gamma[(n-1)/2]) (\sin \theta)^{n-2}$,

\[
\begin{aligned}\int\limits_{\Sigma_{n-1}}e^{-2\pi i r s(x^{\prime}\cdot u^{\prime})}du^{\prime}&=\int\limits_{0}^{\pi}e^{-2\pi i r s\cos\theta}\omega_{n-2}(\sin\theta)^{n-2}d\theta\\&=\omega_{n-2}\int\limits_{-1}^{1}e^{2\pi i r s\xi}(1-\xi^{2})^{(n-3)/2}d\xi\\&=\frac{2\pi^{(n-1)/2}\Gamma[(n-1)/2]\Gamma(\frac{1}{2})}{\Gamma[(n-1)/2](\pi r s)^{(n-2)/2}}J_{(n-2)/2}(2\pi r s)\\&=2\pi(r s)^{-(n-2)/2}J_{(n-2)/2}(2\pi r s).\end{aligned}
\]

We therefore obtain the following result:

\begin{theorem}
\label{thm:c4-3-3}
Suppose $f$ is a radial function in $L^{1}(E_{n})$, $n \geq 2$; thus, $f(x) = f_{0}(|x|)$ for $a.e.x \in E_{n}$. Then the Fourier transform $\hat{f}$ is also radial and has the form $\hat{f}(x) = F_{0}(|x|)$ for all $x \in E_{n}$, where
\[
F_{0}(|x|)= F_{0}(r)=2\pi r^{-[(n-2)/2]}\int_{0}^{\infty}f_{0}(s)J_{(n-2)/2}(2\pi r s)s^{n/2} d s.\footnote{The reason why both here and in Theorem \ref{thm:c4-1-6} we consider the action of the Fourier transform on $L^{1}(E_{n})$, rather than on $L^{2}(E_{n})$, is that in this way all integrals encountered are defined in the $L^{1}$ sense. These results, however, have obvious extensions to $L^{2}(E_{n})$.}
\]

This theorem shows how the Fourier transform acts on the space $\mathfrak{S}_0$. We now turn our attention to the spaces $\mathfrak{S}_k, k \geq 1$. We begin by first examining the action of the Fourier transform on an important class of functions:
\end{theorem}

\begin{theorem}
\label{thm:c4-3-4}
Suppose $f(u) = e^{-\pi|u|^{2}}P(u)$ for $u$ in $E_n$, where $P(u)$ is a solid spherical harmonic of degree $k$, then $\hat{f}(v) = i^{-k}f(v)$ for all $v$ in $E_n$.
\end{theorem}

\textsc{Proof.}\ For $t \in E_n$ fixed we have

\[
\int_{E_{n}}e^{-\pi|u|^{2}}P(u+t)\,du=\int_{0}^{\infty}r^{n-1}e^{-\pi r^{2}}\Biggl\{\int_{\Sigma_{n-1}}P(t+ru^{\prime})\,du^{\prime}\Biggr\}\,dr.
\]

Since $P$ is harmonic it satisfies the mean-value property (Theorem \ref{thm:c2-1-1} of Chapter II) and, thus, $\int_{\Sigma_{n-1}}P(t+ru')du'=\omega_{n-1}P(t)$. Consequently,

\[
\begin{aligned}\int\limits_{E_{n}}e^{-\pi|u|^{2}}P(u+t)du&=P(t)\int\limits_{0}^{\infty}r^{n-1}e^{-\pi r^{2}}\omega_{n-1}dr\\&=P(t)\int\limits_{0}^{\infty}r^{n-1}e^{-\pi r^{2}}\Biggl\{\int\limits_{\Sigma_{n-1}}du^{\prime}\Biggr\}dr\\&=P(t)\int\limits_{E_{n}}e^{-\pi|x|^{2}}dx=P(t).\end{aligned}
\]

$P(t) = P(t_1, \ldots, t_n)$, being a polynomial, has an obvious analytic continuation $P(z) = P(z_1, z_2, \ldots, z_n)$ to all of $C_n$. It then follows from the

last equality that

\[
\int\limits_{E_{n}}e^{-\pi|u|^{2}}P(u+z)d u=P(z)
\]

for all $z = x + iy = (x_1 + iy_1, x_2 + iy_2, \ldots, x_n + iy_n) \in \mathbb{C}_n$. In particular for $z = -iv$ we have

\[
\int\limits_{E_{n}}e^{-\pi|u|^{2}}P(u-i v)d u=P(-i v)=(-i)^{k}P(v),
\]

the last equality being a consequence of the homogeneity of $P$. But an n-fold application of Cauchy's integral theorem shows that

\[
\int_{E_{n}}e^{-\pi(u+i v)\cdot(u+i v)}P(u)\;d u=\int_{E_{n}}e^{-\pi u\cdot u}P(u-i v)\;d u=(-i)^{k}P(v),
\]

where $(u + iv) \cdot (u + iv) = \sum_{j=1}^{n} (u_j + iv_j)^2$. If we now multiply the left and right sides by $e^{-\pi|v|^2}$ we obtain the desired equality $\hat{f}(v) = (-i)^k f(v) = i^{-k} f(v)$.

Theorem \ref{thm:c4-3-4} enables us to study the behavior of the Fourier transform on a still larger class of functions. For $\alpha > 0$ let $g(x) = e^{-\pi |\alpha x|^2} P(x) = \alpha^{-k} f(\alpha x)$ for $x$ in $E_n$. If $\delta_\alpha$ denotes dilation by $\alpha$ it then follows from \eqref{eq:c1-1-6} of Chapter I, Theorem \ref{thm:c4-3-4} and the homogeneity of $P(x)$ that

\begin{equation}\label{eq:c4-3-5}
\hat{g}(x)=\alpha^{-k}(\delta_{\alpha}f)^{\wedge}(x)=\alpha^{-k}\alpha^{-n}\hat{f}(\alpha^{-1}x)=i^{-k}\alpha^{-n-2k}e^{-\pi|x|^{2}/\alpha^{2}}P(x).
\end{equation}

On the other hand, if $h$ is the radial function whose values are $h(x) = e^{-\pi |\alpha x|^2}$ for $x \in E_{n+2k}$, then an application of Theorem \ref{thm:c1-1-13} of Chapter I gives us

\begin{equation}\label{eq:c4-3-6}
\begin{array}{r}{\hat{h}(x)\;=\;\alpha^{-n-2k}e^{-\pi|x|^{2}/\alpha^{2}}.}\end{array}
\end{equation}

We can interpret these two identities in the following way: Let $\mathcal{D}_m$ be the Hilbert space of functions $\varphi$ on $(0, \infty)$ satisfying

\[
\|\varphi\|=\left(\int\limits_{0}^{\infty}|\varphi(r)|^{2}r^{m-1}d r\right)^{1/2}<\infty.
\]

If $m = n + 2k$ and $P(x)$ is a nonzero solid spherical harmonic of degree $k$ we consider the function $g$ whose values are $g(x) = \varphi(|x|)P(x)$ for $x$ in $E_n$. Then, if $\|P\|_{\Sigma} = (\int_{\Sigma_{n-1}}|P(x')|^2 \, dx')^{1/2}$,

\[
\|g\|_{2}^{2}=\int_{E_{n}}|\varphi(|x|)P(x)|^{2}d x=\left(\int_{0}^{\infty}|\varphi(r)|^{2}r^{n+2k-1}d r\right)\|P\|_{\Sigma}^{2}<\infty.
\]

Thus, $g \in L^{2}(E_{n})$ and, by the Plancherel theorem, $\hat{g} \in L^{2}(E_{n})$ with $\|\hat{g}\|_{2} = \|g\|_{2}$. It follows from \eqref{eq:c4-2-19} and this equality between the norms of $g$ and $\hat{g}$ that $\hat{g}(x) = \psi(|x|)P(x)$ for almost every $x$ in $E_{n}$ and that $\|\psi\| = \|\varphi\|$. We therefore obtain a bounded linear operator, $T_{k}^{n}$, on $\mathcal{D}_{n+2k}$ by letting $T_{k}^{n}\varphi = \psi$ (in fact, $T_{k}^{n}$ is an isometry). In a similar way, by considering the radial function $h$ defined by letting $h(x) = \varphi(|x|)$ for $x$ in $E_{n+2k}$ we obtain another bounded linear operator, $T_{0}^{n+2k}$, on $\mathcal{D}_{n+2k}$. In order to define this operator we first observe that it follows from Corollary \ref{cor:c4-1-2} that $\hat{h}$ is also radial, say $\hat{h}(x) = \theta(|x|)$ for almost every $x$ in $E_{n+2k}$. Putting $T_{0}^{n+2k}\varphi = \theta$ we obtain the desired bounded linear operator on $\mathcal{D}_{n+2k}$ (again, it is an immediate consequence of the Plancherel theorem that $T_{0}^{n+2k}$ is an isometry).

Equalities \eqref{eq:c4-3-5} and \eqref{eq:c4-3-6} show that $T_0^n + 2k\varphi = i^k T_k^n\varphi$ whenever $\varphi(r) = e^{-\varepsilon r^2}$, $\varepsilon > 0$. Thus, $T_0^n + 2k$ and $i^k T_k^n$ agree when they are restricted to the space $\mathcal{W}$ of finite linear combinations of all such functions $\varphi$ obtained by varying $\varepsilon$ throughout the positive real numbers. But it is not hard to show that $\mathcal{W}$ is dense in the Hilbert space $\mathcal{D}_{n+2k}$. If this were not the case, we could find a $b \in \mathcal{D}_{n+2k}$, not equal to 0 almost everywhere, such that $\int_0^\infty \varphi(r) b(r) r^{n+2k-1} dr = 0$ for all $\varphi$ in $\mathcal{W}$. In particular,

\refstepcounter{equation}% the book prints no (3.7)
\begin{equation}\label{eq:c4-3-8}
\int\limits_{0}^{\infty}e^{-\varepsilon r^{2}}b(r)r^{n+2k-1}dr=0
\end{equation}

for all $\varepsilon > 0$. Let $\Phi$ be the function whose value at $s \geq 0$ is $\Phi(s) = \int_0^s e^{-r^2} b(r) r^{n+2k-1} dr$. Then, putting $\varepsilon = (m+1)$, where $m$ is a positive integer, and integrating the integrand in \eqref{eq:c4-3-8} by parts we have

\[
0=\int\limits_{0}^{\infty}e^{-m r^{2}}\Phi^{\prime}(r)\,d r=2m\int\limits_{0}^{\infty}e^{-m r^{2}}\Phi(r)r\,d r.
\]

By the change of variable $u = e^{-r^2}$ this equality is equivalent to

\[
0=\int\limits_{0}^{1}u^{m-1}\Phi\biggl(\sqrt{\log\frac{1}{u}}\biggr)\;d u\qquad\text{for}\;m=1,2,3,\ldots.
\]

Since the polynomials are uniformly dense in the space of continuous functions on the closed interval [0,1], this can only be the case when $\Phi(\sqrt{\log 1/u}) = 0$ for all $u$ in [0,1]. Thus $\Phi'(r) = e^{-r^2}b(r)r^{n+2k-1} = 0$

for almost every $r$ in $(0, \infty)$, contradicting the hypothesis that $b(r)$ is not equal to 0 almost everywhere.

Since the operators $T_0^{n+2k}$ and $i^k T_k^n$ are bounded and agree on the dense subspace $\mathcal{W}$ they must be equal. Thus, we have shown

\begin{equation}\label{eq:c4-3-9}
T_{0}^{n+2k}\varphi=i^{k}T_{k}^{n}\varphi\qquad\text{ for all }\varphi\in\mathcal{D}_{n+2k}.
\end{equation}

In Theorem \ref{thm:c4-3-3} we expressed the Fourier transform of a radial function in terms of an integral involving a Bessel function. Equalities \eqref{eq:c4-2-19} and \eqref{eq:c4-3-9} show that the Fourier transform of a function in $\mathfrak{D}_k$ of the form $f_0(|x|)P(x)$, where $P(x)$ is a solid harmonic of degree $k$ on $E_n$, can be expressed in terms of the Fourier transform of the radial function whose values are $h(y) = f_0(|y|)$ for $y \in E_{n+2k}$. Combining these results we obtain the following theorem (which has an obvious extension to all of $L^1(E_n)$ or $L^2(E_n)$):

\begin{theorem}
\label{thm:c4-3-10}
Suppose $n \geq 2$ and $f \in L^2(E_n) \cap L^1(E_n)$ has the form $f(x) = f_0(|x|)P(x)$, where $P(x)$ is a solid spherical harmonic of degree $k$, then $\hat{f}$ has the form $\hat{f}(x) = F_0(|x|)P(x)$ where
\[
F_{0}(r)\;=\;2\pi i^{-k}r^{-[(n+2k-2)/2]}\intop_{0}^{\infty}f_{0}(s)J_{(n+2k-2)/2}(2\pi r s)s^{(n+2k)/2}\;d s.
\]

\end{theorem}

This theorem, together with the fact that the spaces $\mathfrak{G}_k$ are spanned by functions of the form $f_0(|x|)P(x)$, gives the promised description of the action of the Fourier transform on $\mathfrak{G}_k$.

The identity describing the action of the Fourier transform on $\mathfrak{V}_k$ naturally leads us to the study of the behavior of the Bessel functions $J_m(r)$. Near the origin we have the trivial estimate $J_m(r) \sim cr^m$ as $r \to 0$. The behavior for large r is less obvious and is given in the following lemma.

\begin{lemma}
\label{lem:c4-3-11}
$J_{m}(r) = \sqrt{2/\pi r} \cos (r - \pi m/2 - \pi/4) + O(r^{-3/2})$ as $r \to \infty$. In particular,
\begin{equation}\label{eq:c4-3-12}
J_{m}(r)=O(r^{-1/2})\qquad\text{as}\qquad r\to\infty.\footnote{For later purposes it is useful to observe that the proof given below shows that the bounds arising in the terms $O(r^{-1/2})$ and $O(r^{-3/2})$ are uniform in $m$, as long as $m$ is restricted to any closed, bounded subinterval of $(-\frac{1}{2}, \infty)$.}
\end{equation}

\end{lemma}

\textsc{Proof.}\ In view of the definition of $J_{m}(r)$ for $m > -\frac{1}{2}$, we need to estimate the integral

\[
I=\int\limits_{-1}^{1}e^{i r s}(1-s^{2})^{m-1/2}\,ds.
\]

In order to do this we consider the simply connected region in the complex plane obtained by excluding the rays $(-\infty, -1)$ and $(1, \infty)$. We then

choose the branch of $f(z) = (1 - z^2)^{m-1/2}$ in this region which is real valued and nonnegative on the interval $[-1, 1]$. Integrating $e^{irz}(1 - z^2)^{m-1/2} = e^{ir(x+iy)}(1 - [x + iy]^2)^{m-1/2}$ around the rectangle whose lower side is $[-1, 1]$ and whose height is $a > 0$, we obtain

\[
\begin{aligned}&0=\intop_{0}^{a}e^{i r(1+i y)}(y^{2}-2i y)^{m-1/2}d i y\\
&\quad{}+\intop_{-1}^{1}e^{i r s}(1-s^{2})^{m-1/2}d s\\
&\quad{}+\intop_{a}^{0}e^{i r(i y-1)}(y^{2}+2i y)^{m-1/2}d i y+\varepsilon(a),\end{aligned}
\]

where $\varepsilon(a) \to 0$ as $a \to \infty$. Thus,

\[
\begin{aligned}&-i I=\int\limits_{0}^{\infty}e^{i r(i y-1)}(y^{2}+2i y)^{m-1/2}d y\\
&\qquad-\int\limits_{0}^{\infty}e^{i r(1+i y)}(y^{2}-2i y)^{m-1/2}d y\\
&\qquad=I_{1}-I_{2}.\end{aligned}
\]

Now

\[
(y^{2}\pm2i y)^{m-1/2}=\left\{\begin{matrix}{y^{m-1/2}(\pm2i)^{m-1/2}+O(y^{m+1/2}),\text{for }0\leq y\leq1}\\ {y^{m-1/2}(\pm2i)^{m-1/2}+O(y^{2m-1}),\text{for }1\leq y<\infty.}\\ \end{matrix}\right.
\]

Therefore,

\[
\begin{aligned}{I_{1}= e^{-i r}\int_{0}^{\infty}e^{-r y}(2i)^{m-1/2}y^{m-1/2}\,dy+O\bigg(}&{{}\int_{0}^{1}e^{-r y}y^{m+1/2}\,dy\bigg)}\\ {}&{{}+O\bigg(\int_{1}^{\infty}e^{-r y}y^{2m-1}\,dy\bigg).}\\ \end{aligned}
\]

The first term is $(2i)^{m-1/2}\Gamma(m+\frac{1}{2})(e^{-ir}/r^{m+1/2})$. The second is $O(r^{-m-3/2})$ and the third is $O(e^{-r})$ as $r\to\infty$. Similarly, $I_{2}=(-2i)^{m-1/2}\Gamma(m+\frac{1}{2})\times(e^{ir}/r^{m+1/2})+O(r^{-m-3/2})$, as $r\to\infty$. Since $J_{m}(r)=[(r/2)^{m}/\Gamma(m+\frac{1}{2})\Gamma(\frac{1}{2})]\times i(I_{1}-I_{2})$, we obtain the desired result by collecting terms.

\section{Some Applications}

Many important operators in analysis are convolution operators. That is, they map functions $f$, ranging throughout some function space, into convolutions $K* f$, where $K$ is a fixed function (or, more generally, a tempered distribution). We have encountered several such operators. One example is given by the Poisson integral, where $K(x) = K_y(x) = c_ny/(|x|^2 + y^2)^{(n+1)/2}$ for some $y > 0$ (see Section 1 in Chapter I). If $f$

belongs to $L^{1}(E_{n})$ (or to $L^{2}(E_{n})$) it follows from Theorems \ref{thm:c1-1-4} and 1.14 of Chapter I that

\[
(K* f)\hat{}^{\wedge}(x)=e^{-2\pi y\vert x\vert}\hat{f}(x)
\]

for $x$ in $E_n$ (or a.e. $x$ in $E_n$). We see, therefore, that if we follow this convolution operator with the Fourier transform we obtain a transformation of a very simple form. For example, if $f \in L^2(E_n)$ it is an immediate consequence of the Plancherel theorem and the convergence of $e^{-2\pi y|x|}$ to 1, as $y$ tends to 0, that $\lim_{y \to 0} \|f - (K_y * f)\|_2 = 0$. This is a particularly simple example of a general situation: After taking the Fourier transform of a convolution operator we obtain a multiplication operator, many of whose important properties are easily derived. The Poisson kernel is a radial function. The other kernels we shall consider are either radial or are products of radial functions with spherical harmonics. In view of Theorem \ref{thm:c4-3-10} it is not unreasonable to expect that the Fourier transform of such a kernel has the same form (even when the kernel is not an $L^1$ or $L^2$ function).

Our first application concerns a class of kernels that give rise to convolution transforms that are important in Fourier analysis and the theory of partial differential equations. These are the functions of the form $P(x)/|x|^\beta$, where $P$ is a solid spherical harmonic of degree $k$ and $\beta$ is a complex number. If $K$ is such a function and $\text{Re}(\beta) < n + k$, then it is locally integrable and it defines a tempered distribution in the sense described in example (3) of the third section in Chapter I. Thus, its Fourier transform $\hat{K}$ is well-defined and is also a tempered distribution. The following result asserts that $K$ is a tempered function for a certain range of $\beta$ and gives its explicit form.

\begin{theorem}
\label{thm:c4-4-1}
Suppose $\alpha$ is a complex number such that $0 < \text{Re}(\alpha) < n$ and $P(x)$ is a harmonic polynomial on $E_n$ that is homogeneous of degree $k$. If $K(x) = P(x)/|x|^{n+k-\alpha}$ then $\hat{K}(t) = \gamma P(t)/|t|^{k+\alpha}$, where
\[
\gamma=\gamma_{k,\alpha}=i^{-k}\pi^{(n/2)-\alpha}\Gamma\bigg(\frac{k+\alpha}{2}\bigg)\bigg/\Gamma\bigg(\frac{n+k-\alpha}{2}\bigg).
\]

\end{theorem}

\textsc{Proof.}\ We define $K_1$ by letting $K_1(x) = K(x)$ if $|x| \leq 1$ and $K_1(x) = 0$ if $|x| > 1$; we then define $K_2$ to be the difference $K - K_1$. $K_1$ then belongs to $L^1(E_n)$ and, if we impose the additional restriction that $\text{Re}(\alpha) < n/2$, we can then conclude that $K_2$ belongs to $L^2(E_n)$. Thus it follows from the definition of the Fourier transform of a distribution (see \eqref{eq:c1-3-15}, Chapter I) that $\hat{K}$ is the function $\hat{K}_1 + \hat{K}_2$. By approximating $K_1$ and $K_2$ with functions in $L^1(E_n) \cap L^2(E_n)$ and using \eqref{eq:c4-2-19} we see that $\hat{K}_1$ and $\hat{K}_2$ must have the forms $\hat{K}_1(x) = f_1(|x|)P(x)$ and $\hat{K}_2(x) = f_2(|x|)P(x)$. Thus, $\hat{K}(x) = f(|x|)P(x)$, where $f = f_1 + f_2$. If $K$ were in $L^1(E_n) \cap L^2(E_n)$ it would be an

immediate consequence of Theorem \ref{thm:c4-3-10} that $f$ is a homogeneous function of degree $-(k+\alpha)$. Unfortunately this is not the case; thus, we cannot apply that theorem directly in order to obtain this homogeneity property. Although we could use Theorem \ref{thm:c4-3-10} and an approximation argument to show that for a.e. $r>0$

\begin{equation}\label{eq:c4-4-2}
f(\delta r)=\delta^{-k-\alpha}f(r)
\end{equation}

whenever $\delta > 0$, it is easier to do this by exploiting the relation between the Fourier transform and the action of dilations on $E_n$.

Let $\varphi$ be a testing function (a member of $\mathcal{S}$). By the multiplication formula (in fact, by the definition of the Fourier transform of tempered distributions) we have

\[
\int_{E_{n}}K(x)\hat{\varphi}(x)\;d x=\int_{E_{n}}\hat{K}(x)\varphi(x)\;d x.
\]

Putting $x = \delta u$ on the left and $x = \delta^{-1}v$ on the right, this equality becomes

\[
\delta^{n}\int_{E_{n}}K(\delta u)\hat{\varphi}(\delta u)\;d u=\;\delta^{-n}\int_{E_{n}}\hat{K}(\delta^{-1}v)\varphi(\delta^{-1}v)\;d v.
\]

Since $K(\delta u) = \delta^{\alpha-n}K(u)$

\[
\delta^{\alpha}\int_{E_{n}}K(u)\hat{\varphi}(\delta u)\;d u\;=\;\delta^{-n}\int_{E_{n}}\hat{K}(\delta^{-1}v)\varphi(\delta^{-1}v)\;d v.
\]

On the other hand, the Fourier transform of the function $\psi$ whose values are $\varphi(\delta^{-1}v)$ is the function defined by $\hat{\psi}(u) = \delta^n \hat{\varphi}(\delta u)$ for $u \in E_n$. Thus,

\[
\begin{aligned}{\delta^{-\alpha}\intop_{E_{n}}\hat{K}(\delta^{-1}v)\varphi(\delta^{-1}v)\;d v}&{{}=\;\delta^{n}\intop_{E_{n}}K(u)\hat{\varphi}(\delta u)\;d u=\intop_{E_{n}}K(u)\;\hat{\psi}(u)\;d u}\\ {}&{{}=\;\intop_{E_{n}}\hat{K}(u)\psi(u)\;d u=\intop_{E_{n}}\hat{K}(v)\varphi(\delta^{-1}v)\;d v}\\ \end{aligned}
\]

for all testing functions $\varphi$. Consequently, $\hat{K}(v) = \delta^{-\alpha}\hat{K}(\delta^{-1}v)$ for a.e. $v \in E_n$ and \eqref{eq:c4-4-2} follows immediately. Thus,

\[
F(\xi)=\int\limits_{0}^{\xi}f(r)d r=\int\limits_{0}^{1}f(\xi s)\xi d s=F(1)\xi^{1-k-\alpha}
\]

for all $\xi > 0$ and, consequently, $f(r) = F'(r) = \gamma r^{-k-\alpha}$. This shows that $\hat{K}(t) = \gamma P(t)/|t|^{k+\alpha}$ for a.e. $t \in E_n$, provided $0 < \text{Re}(\alpha) < n/2$. In order to evaluate $\gamma$, we let $\varphi(x) = e^{-\pi|x|^2}P(x)$. By Theorem \ref{thm:c4-3-4} we then have

$\hat{\varphi}(t) = i^{-k}\varphi(t)$; hence,

\[
i^{-k}\int\limits_{E_{n}}e^{-\pi|t|^{2}}P(t)\{P(t)/|t|^{k+n-\alpha}\}d t=\gamma\int\limits_{E_{n}}e^{-\pi|t|^{2}}P(t)\{P(t)/|t|^{k+\alpha}\}d t.
\]

Expressing these two integrals in polar coordinates and then cancelling the common term $\int_{\Sigma_{n-1}} [P(t')]^2 dt'$, we obtain

\[
i^{-k}\int\limits_{0}^{\infty}r^{\alpha+k-1}e^{-\pi r^{2}}d r=\gamma\int\limits_{0}^{\infty}r^{k+n-\alpha-1}e^{-\pi r^{2}}d r.
\]

The value of $\gamma$ is easily deduced from this equality and the relation

\[
\Gamma\biggl(\frac{\beta+1}{2}\biggr)\pi^{-(\beta+1)/2}=2\intop_{0}^{\infty}e^{-\pi t^{2}}t^{\beta}d t.
\]

Finally, we must extend our result to the entire range $0 < \text{Re}(\alpha) < n$. We have established the equality

\begin{equation}\label{eq:c4-4-3}
\int\limits_{E_{n}}\left\{P(x)/|x|^{k+n-\alpha}\right\}\hat{\varphi}(x)d x=\gamma_{k,\alpha}\int\limits_{E_{n}}\left\{P(x)/|x|^{k+\alpha}\right\}\varphi(x)d x
\end{equation}

for $0 < \text{Re}(\alpha) < n/2$ when $\varphi$ is a testing function. But each of these two integrals defines an analytic function of $\alpha$ for $0 < \text{Re}(\alpha) < n$. Since $\gamma_{\kappa,\alpha}$ also defines such an analytic function, equality \eqref{eq:c4-4-3} must hold for $0 < \text{Re}(\alpha) < n$. The theorem now follows immediately from the definition of the Fourier transform of a distribution.

Kernels of the form $K(x) = P(x)/|x|^{n+k}$, $k \geq 1$, give rise to an important class of operators in analysis. \footnote{These operators are special cases of the singular integral operators that will be discussed in Chapter VI.} In view of Theorem \ref{thm:c4-4-1} it is natural to try to justify taking limits, as $\alpha$ tends to 0, and claim that

\[
\hat{K}(t)\;=\;i^{-k}\pi^{n/2}\biggl\{\Gamma\biggl(\frac{k}{2}\biggl)\biggr/\Gamma\biggl(\frac{n+k}{2}\biggl)\biggr\}P(t)/|t|^{k}.
\]

The first difficulty we encounter is that K is not integrable in a neighborhood of the origin. The function K, however, can be used in order to define a tempered distribution L. This is done by letting

\begin{equation}\label{eq:c4-4-4}
L(\varphi)=\lim_{\varepsilon\to0}\int\limits_{0<\varepsilon\leq|t|}K(t)\varphi(t)dt
\end{equation}

for each testing function $\varphi$. To see that this limit exists and defines a continuous linear functional on the space $\mathcal{S}$ of testing functions we first observe that it suffices to show that this is true for $\lim\limits_{\varepsilon\to0}\int_{\varepsilon\leq|t|\leq1}K(t)\varphi(t)dt$ (for the function $K_1$, defined by letting $K_1(t)=0$ if $|t|\leq1$ and $K_1(t)=K(t)$ if $|t|>1$, is a tempered function). Since $k\geq1$ we must have

\[
\int_{\varepsilon\leq|t|\leq1}K(t)dt=\int_{\varepsilon}^{1}r^{-1}\left\{\int_{\Sigma_{n-1}}P(t^{\prime})dt^{\prime}\right\}dr=0.
\]

Thus,

\[
\int_{\varepsilon\leq|t|\leq1}K(t)\;\varphi(t)\;{d}t=\int_{\varepsilon\leq|t|\leq1}K(t)[\varphi(t)\;-\;\varphi(0)]\;{d}t.
\]

But

\[
\left|\varphi(t)-\varphi(0)\right|\leq\left\{\sup_{|s|\leq1}\left|(\nabla\varphi)(s)\right|\right\}\big|t\big|.
\]

Thus, if

\[
t=\left|t\right|t^{\prime},\left|K(t)[\varphi(t)\right.-\varphi(0)]\left|\right|\leq\left\{\sup_{s\in E_{n}}\left|(\nabla\varphi)(s)\right|\right\}\left|P(t^{\prime})\right|/\left|t\right|^{n-1};
\]

therefore, $K(t)[\varphi(t) - \varphi(0)]$ is locally integrable and

\[
\begin{aligned}\left|\lim_{\varepsilon\to0}\int\limits_{\varepsilon\leq|t|\leq1}K(t)\varphi(t)dt\right|&=\left|\int\limits_{|t|\leq1}K(t)[\varphi(t)-\varphi(0)]dt\right|\\&\leq\left\{\sup_{s\in E_{n}}\left|(\nabla\varphi)(s)\right|\right\}_{|t|\leq1}\int\left\{|P(t^{\prime})|/|t|^{n-1}\right\}dt.\end{aligned}
\]

The fact that L is a tempered distribution\index{principal-value distribution@principal-value distribution} now follows from Theorem \ref{thm:c1-3-11} of Chapter I.

Distributions obtained by the limiting process \eqref{eq:c4-4-4} are called principal value distributions; the notation usually used to indicate that L is such a linear functional, defined in terms of the function K, is

\[
{\cal L}(\varphi)=\mathrm{P.V.}\int_{E_{n}}{\int}K(t)\varphi(t)d t
\]

for $\varphi \in \mathcal{S}$. In such a case we shall denote the distribution also by $K$.

Since $K(t) = P(t)/|t|^{n+k}$, as we have just seen, defines a (principal value) distribution, the Fourier transform $\hat{K}$ is defined as a tempered distribution. The following theorem justifies the passage to the limit, as $\alpha$ approaches 0, that was discussed above:

\begin{theorem}
\label{thm:c4-4-5}
Suppose $P(x)$ is a harmonic polynomial on $E_n$ that is homogeneous of degree $k \geq 1$. If $K(x) = P(x)/|x|^{n+k}$ then
\[
\hat{K}(t)=\biggl\lbrace i^{-k}\pi^{n/2}\Gamma\Bigl(\frac{k}{2}\Bigr)\biggl/{\Gamma\Bigl(\frac{n+k}{2}\Bigr)}\biggr\rbrace P(t)/|t|^{k}=\gamma_{k,0}P(t)/|t|^{k}.
\]

\end{theorem}

\textsc{Proof.}\ We must show that

\[
\mathrm{(i)}\quad\gamma_{k,0}\int\limits_{E_{n}}\{P(t)/|t|^{k}\}\varphi(t)\;d t=\mathrm{P.V.}\int\limits_{E_{n}}\{P(t)/|t|^{n+k}\}\hat{\varphi}(t)\;d t
\]

for all testing functions $\varphi$. Theorem \ref{thm:c4-4-1} asserts that

\[
\text{(ii)}\quad\gamma_{k,\alpha}\int\limits_{E_{n}}\{P(t)/|t|^{k+\alpha}\}\varphi(t)\;d t=\int\limits_{E_{n}}\{P(t)/|t|^{n+k-\alpha}\}\hat{\varphi}(t)\;d t
\]

when $0 < \alpha < n$. As $\alpha \to 0$ the left side of (ii) approaches the left side of (i); therefore, it suffices to show that

\[
\lim_{\alpha\to0}\int_{E_{n}}\left\{P(t)/|t|^{n+k-\alpha}\right\}\hat{\varphi}(t)\;d t=\mathrm{P.V.}\;\int_{E_{n}}\left\{P(t)/|t|^{n+k}\right\}\hat{\varphi}(t)\;d t.
\]

Since the integrals of $P(t)$ over the surfaces of spheres centered at the origin are zero and, as was observed above, $K(t)[\hat{\varphi}(t) - \hat{\varphi}(0)]$ is locally integrable we have, by the dominated convergence theorem,

\[
\begin{aligned}&\lim_{\alpha\to0}\int\limits_{E_{n}}\{P(t)/|t|^{n+k-\alpha}\}\hat{\varphi}(t)dt\\
&=\lim_{\alpha\to0}\int\limits_{|t|\leq1}\{P(t)/|t|^{n+k-\alpha}\}\{\hat{\varphi}(t)-\hat{\varphi}(0)\}dt+\int\limits_{|t|>1}K(t)\hat{\varphi}(t)dt\\
&=\int\limits_{|t|\leq1}K(t)[\hat{\varphi}(t)-\hat{\varphi}(0)]dt+\int\limits_{|t|>1}K(t)\hat{\varphi}(t)dt.\end{aligned}
\]

But it was shown above that the last expression is the desired limit $P.V. \int_{E_n} \{P(t)/|t|^{n+k}\}\hat{\varphi}(t) \, dt$ and the theorem is proved.

If we let $Y^{(k)}(x') = P^{(k)}(x)/|x|^k$, where $x \neq 0$, $x' = x/|x|$ and $P^{(k)}$ is a harmonic polynomial of degree $k \geq 1$ (i.e., $Y^{(k)}$ is a spherical harmonic of degree $k$), we put $\Omega(x') = \sum_{k=1}^{m} Y^{(k)}(x')$. It then follows from Theorem \ref{thm:c4-4-5} that $K(x) = \Omega(x')/|x|^n$ defines a principal value distribution whose Fourier transform is the function

\[
\begin{array}{r}{\hat{K}(x)=\displaystyle\sum_{k=1}^{m}\gamma_{k,0}Y^{(k)}(x^{\prime}).}\end{array}
\]

In view of Corollary \ref{cor:c4-2-3} it is natural to consider, more generally, such kernels $K$ obtained from those functions $\Omega$ in $L^2(\Sigma_{n-1})$ that lie in the closed subspace of $L^2(\Sigma_{n-1})$ generated by the spherical harmonics $Y^{(k)}$ for $k=1,2,3,\ldots$. This is precisely the space of those $L^2$ functions satisfying

\begin{equation}\label{eq:c4-4-6}
\int_{\Sigma_{n-1}}\Omega(x^{\prime})\,dx^{\prime}=0.
\end{equation}

If we let $K(x) = \Omega(x')/|x|^n$, a simple extension of the argument that followed \eqref{eq:c4-4-4} shows that

\[
\lim_{\varepsilon\to0}\int\limits_{0<\varepsilon\leq|t|}K(t)\varphi(t)dt,
\]

where $\varphi$ is a testing function, defines a principal value distribution that will also be denoted by K. We then have the following generalization of Theorem \ref{thm:c4-4-5}:

\begin{theorem}
\label{thm:c4-4-7}
Suppose $\Omega$ is a function in $L^2(\Sigma_{n-1})$ satisfying \eqref{eq:c4-4-6}. Then there exist unique spherical harmonics $Y^{(k)}$, of degree k, such that
\begin{equation}\label{eq:c4-4-8}
\Omega=\sum_{k=1}^{\infty}Y^{(k)},
\end{equation}

the series being convergent in the norm of $L^{2}(\Sigma_{n-1})$. The function whose values at $x \neq 0$ are $K(x) = \Omega(x')/|x|^{n}$ defines a principal value distribution whose Fourier transform $\hat{K}$ is a function that is homogeneous of degree zero; thus, $\hat{K}(x) = \Omega_0(x') = \Omega_0(x/|x|)$ for $x \neq 0$. Moreover, $\Omega_0 = \sum_{k=1}^{\infty} Y_0^{(k)}$, where $Y_0^{(k)} = \gamma_{k,0} Y^{(k)}$ and the series is convergent in the norm of $L^{2}(\Sigma_{n-1})$. Furthermore,
\begin{equation}\label{eq:c4-4-9}
\sum_{k=1}^{\infty}k^{n}\|Y_{0}^{(k)}\|^{2}=\sum_{k=1}^{\infty}k^{n}\int\limits_{\Sigma_{n-1}}\mid Y_{0}^{(k)}(x^{\prime})\mid^{2}d x^{\prime}<\infty.
\end{equation}

\end{theorem}

Conversely, any function, $\Omega_0$, that is homogeneous of degree 0 and whose restriction to $\Sigma_{n-1}$ is square integrable, satisfies $\int_{\Sigma_{n-1}} \Omega_0(x') \, dx' = 0$ and has a spherical harmonic development satisfying \eqref{eq:c4-4-9} arises as the Fourier transform of a principal value distribution having the form $K(x) = \Omega(x') / |x|^n$ where $\Omega$ belongs to $L^2(\Sigma_{n-1})$ and satisfies \eqref{eq:c4-4-6}.

\textsc{Proof.}\ The representation \eqref{eq:c4-4-8} is a consequence of Corollary \ref{cor:c4-2-3} which, together with the orthogonality relations Corollary \ref{cor:c4-2-4}, imply that

\[
\sum_{k=1}^{\infty}\left\|Y^{(k)}\right\|^{2}<\infty.
\]

From Theorem \ref{thm:c4-4-5} we see that $\gamma_{k,0} \approx k^{-n/2}$ as $k$ tends to $\infty$; therefore, inequality \eqref{eq:c4-4-9} must hold. In order to see that $\tilde{K}$ is obtained from the function $\Omega_0$ in the manner described above we consider the kernel $K^{(m)}(x) = \Omega^{(m)}(x') / |x|^n = \sum_{k=1}^m Y^{(k)}(x') / |x|^n$. From Theorem \ref{thm:c4-4-5} we see that the Fourier transform of $K^{(m)}$ (considered as a principal value distribution) is the function, homogeneous of degree 0, whose restriction to $\Sigma_{n-1}$ is $\sum_{k=1}^m Y_0^{(k)} = \sum_{k=1}^m \gamma_{k,0} Y^{(k)}$. Thus, if $\varphi$ is a testing function, it follows from the definition of the Fourier transform of tempered distributions that

\[
\int_{E_{n}}\hat{K}^{(m)}(x)\varphi(x)\;d x=\mathrm{P.V.}\int_{E_{n}}K^{(m)}(x)\hat{\varphi}(x)\;d x.
\]

Since $\sum_{k=1}^{m} Y_0^{(k)} \to \sum_{k=1}^{\infty} Y_0^{(k)} = \Omega_0$ in $L^2(\Sigma_{n-1})$ as $m \to \infty$ (this is an immediate consequence of \eqref{eq:c4-4-9}), it follows that

\[
\lim_{m\to\infty}\int\limits_{E_{n}}\hat{K}^{(m)}(x)\varphi(x)\,dx=\int\limits_{E_{n}}\Omega_{0}(x/|x|)\varphi(x)\,dx.
\]

On the other hand, the same argument we gave after \eqref{eq:c4-4-4} shows that

\begin{equation}\label{eq:c4-4-10}
\begin{aligned}\mathrm{P.V.}\int_{E_{n}}&K^{(m)}(x)\hat{\varphi}(x)dx\\&=\int_{|x|\leq1}K^{(m)}(x)[\hat{\varphi}(x)-\hat{\varphi}(0)]dx+\int_{|x|>1}K^{(m)}(x)\hat{\varphi}(x)dx.\end{aligned}
\end{equation}

But the $L^2$ convergence $\sum_{k=1}^m Y^{(k)} \to \Omega$, as $m \to \infty$, implies that the right side of \eqref{eq:c4-4-10} converges to

\[
\int\limits_{|x|\leq1}K(x)[\hat{\varphi}(x)-\hat{\varphi}(0)]\;{d}x\;+\int\limits_{|x|>1}K(x)\hat{\varphi}(x)\;{d}x=\mathrm{P.V.}\int\limits_{E_{n}}K(x)\hat{\varphi}(x)\;{d}x.
\]

Thus,

\[
\int_{E_{n}}\Omega_{0}(x/|x|)\varphi(x)\;d x=\mathrm{P.V.}\int_{E_{n}}K(x)\hat{\varphi}(x)\;d x
\]

for all testing functions $\varphi$ and it follows from the definition of the Fourier transform of the distribution $K$ and $\hat{K}$ is the function whose values are $\Omega_0(x/|x|)$ when $x \neq 0$.

In order to show the last part we first observe that if $\Omega_0$ has a spherical harmonic development $\sum_{k=1}^{\infty} Y_0^{(k)}$ satisfying \eqref{eq:c4-4-9} then $\sum_{k=1}^{\infty} Y^{(k)}$, where $Y^{(k)} = \gamma_{k,0}^{-1} Y_0^{(k)}$, converges in $L^2(\Sigma_{n-1})$ to a function $\Omega$ having integral zero over $\Sigma_{n-1}$. By applying the first part of our theorem to this function $\Omega$ we then obtain the desired converse.

More generally, the argument we gave showing that

\[
\begin{aligned}\mathrm{P.V.}\int_{E_{n}}K(x)\varphi(x)dx&=\lim_{\varepsilon\to0}\int\limits_{|x|\geq\varepsilon>0}K(x)\varphi(x)dx\\&=\int\limits_{|x|\leq1}K(x)[\varphi(x)-\varphi(0)]dx+\int\limits_{|x|>1}K(x)\varphi(x)dx,\end{aligned}
\]

when $\varphi$ is a testing function and $K(x) = \Omega(x')/|x|^n$, depended only on the integrability of $\Omega$ on $\Sigma_{n-1}$ and property \eqref{eq:c4-4-6}. Thus, $K$ is a principal value distribution in this case as well; the following theorem gives an explicit form for its Fourier transform $\hat{K}$:

\begin{theorem}
\label{thm:c4-4-11}
Suppose $\Omega \in L^{1}(\Sigma_{n-1})$ satisfies \eqref{eq:c4-4-6}. Then the function whose values at $x \neq 0$ are $K(x) = \Omega(x')/|x|^{n}$ defines a principal value distribution whose Fourier transform $\hat{K}$ is a function that is homogeneous of degree zero; thus, $\hat{K}(x) = \Omega_0(x') = \Omega_0(x/|x|)$ for $x \neq 0$. Moreover,
\[
\Omega_{0}(x^{\prime})\;=\;-\;\intop_{\Sigma_{n-1}}\Omega(t^{\prime})\biggl[\frac{i\pi}{2}\text{sgn}\left(x^{\prime}\cdot t^{\prime}\right)\;+\;\log\left|x^{\prime}\cdot t^{\prime}\right|\biggr]\;d t^{\prime}
\]

for $x' \in \Sigma_{n-1}$, where $\text{sgn}$ denotes the signum function on $(-\infty, \infty)$: $\text{sgn}(s) = 1$ if $s > 0$, $\text{sgn}(s) = -1$ if $s < 0$ and $\text{sgn}(0) = 0$.
\end{theorem}

\textsc{Proof.}\ Let $\hat{K}_{\varepsilon}^{N}$ be the Fourier transform of the $L^{1}(E_{n})$ function whose values are $\Omega(x^{\prime})/|x|^{n}$, when $0 < \varepsilon < |x| < N$, and is 0 when $|x|$ is outside this interval.

If we let $x = rx'$ and $t = \rho t'$, with $x'$, $t'$ on $\Sigma_{n-1}$, we have

\[
\begin{aligned}{\hat{K}_{\varepsilon}^{N}(x)}&{{}=\intop_{\varepsilon<|t|<N}e^{-2\pi i x\cdot t}\frac{\Omega(t^{\prime})}{|t|^{n}}d t}\\ {}&{{}=\intop_{\Sigma_{n-1}}\Omega(t^{\prime})\Biggl\{\intop_{\varepsilon}^{N}\frac{e^{-2\pi i r\rho(x^{\prime}\cdot t^{\prime})}}{\rho}d\rho\Biggr\}d t^{\prime}.}\\ \end{aligned}
\]

Since $\Omega$ satisfies \eqref{eq:c4-4-6} this last integral equals

\[
\begin{aligned}&{\intop_{\Sigma_{n-1}}\Omega(t^{\prime})\Bigg\{\intop_{\varepsilon}^{N}\frac{\cos2\pi r\rho(x^{\prime}\cdot t^{\prime})-\cos2\pi r\rho}{\rho}d\rho\Bigg\}d t^{\prime}}\\ {}&{{}-i\intop_{\Sigma_{n-1}}\Omega(t^{\prime})\Bigg\{\intop_{\varepsilon}^{N}\frac{\sin2\pi r\rho(x\cdot^{\prime}t^{\prime})}{\rho}d\rho\Bigg\}d t^{\prime}.}\\ \end{aligned}
\]

But, as is well-known,

\[
\left|\int\limits_{\varepsilon}^{N}\frac{\sin2\pi r\rho(x^{\prime}\cdot t^{\prime})}{\rho}d\rho\right|=\left|\int\limits_{2\pi r\varepsilon(x^{\prime}\cdot t^{\prime})}^{2\pi r N(x^{\prime}\cdot t^{\prime})}\frac{\sin s}{s}ds\right|
\]

is dominated by a constant independent of $\varepsilon$, N, r and $x' \cdot t'$ and

\[
\lim_{\varepsilon\to0,N\to\infty}\;\int_{\varepsilon}^{N}\frac{\sin2\pi r\rho(x^{\prime}\cdot t^{\prime})}{\rho}\;d\rho\;=\frac{\pi}{2}\text{sgn}(x^{\prime}\cdot t^{\prime}).
\]

Moreover,

\[
\begin{aligned}&\int_{\varepsilon}^{N}\frac{\cos 2\pi r\rho(x^{\prime}\cdot t^{\prime})-\cos 2\pi r\rho}{\rho}\,d\rho\\
&=\int_{2\pi r\varepsilon|x^{\prime}\cdot t^{\prime}|}^{2\pi rN|x^{\prime}\cdot t^{\prime}|}\frac{\cos s}{s}\,ds-\int_{2\pi r\varepsilon}^{2\pi rN}\frac{\cos s}{s}\,ds\\
&=\int_{2\pi r\varepsilon|x^{\prime}\cdot t^{\prime}|}^{2\pi r\varepsilon}\frac{\cos s}{s}\,ds-\int_{2\pi rN|x^{\prime}\cdot t^{\prime}|}^{2\pi rN}\frac{\cos s}{s}\,ds.\end{aligned}
\]

Each of these last two integrals has absolute value dominated by $\log\left(1/|x' \cdot t'|\right)$ while an easy calculation shows that

\[
\lim_{\varepsilon\to0}\int_{2\pi r\mid x^{\prime}\cdot t^{\prime}\mid\varepsilon}^{2\pi r\varepsilon}\frac{\cos s}{s}ds=\log\frac{1}{\left|x^{\prime}\cdot t^{\prime}\right|}
\]

and

\[
\lim_{N\to\infty}\int_{2\pi r|x^{\prime}\cdot t^{\prime}|N}^{2\pi r N}\frac{\cos s}{s}d s=0.
\]

We have shown, therefore, that

\[
\lim_{\varepsilon\to0,N\to\infty}\int_{\varepsilon}^{N}\frac{e^{-2\pi i r\rho\left(x^{\prime}\cdot t^{\prime}\right)}}{\rho}d\rho\;=\;-\;\log\left|x^{\prime}\cdot t^{\prime}\right|\;-\;i\frac{\pi}{2}\text{sgn}\left(x^{\prime}\cdot t^{\prime}\right)
\]

and that this convergence is dominated by a constant (independent of $\epsilon$, $N$, $r$, and $x' \cdot t'$) times $1 + \log(1/|x' \cdot t'|)$.

It is an immediate consequence of Fubini's theorem that the integral $\int_{\Sigma_{n-1}}|\Omega(t')|[1 + \log (1/|x'\cdot t'|)]\,dt'$ is finite for almost every $x'$ in $\Sigma_{n-1}$ and represents an integrable function on $\Sigma_{n-1}$. Thus, by the Lebesgue dominated convergence theorem

\[
\lim_{\varepsilon\to0,N\to\infty}\hat{K}_{\varepsilon}^{N}(x^{\prime})\;=\;-\intop_{\Sigma_{n-1}}\Omega(t^{\prime})\biggl[\frac{i\pi}{2}\text{sgn}\left(x^{\prime}\cdot t^{\prime}\right)\;+\;\log\left|x^{\prime}\cdot t^{\prime}\right|\biggr]\;{d}t^{\prime}.
\]

On the other hand, by the multiplication formula,

\[
\int_{E_{n}}\hat{K}_{\varepsilon}^{N}(x^{\prime})\varphi(x)d x=\int_{\varepsilon<|x|<N}\frac{\Omega(x^{\prime})}{|x|^{n}}\,\hat{\varphi}(x)d x
\]

whenever $\varphi$ is a testing function. Theorem \ref{thm:c4-4-11} now follows easily by passing to the limit as $\varepsilon \to 0$ and $N \to \infty$ and using the definition of the Fourier transform of the principal value distribution $K$.

\begin{corollary}
\label{cor:c4-4-12}
If $\Omega \in L^2(\Sigma_{n-1})$ then the function $\Omega_0$ is continuous.
\end{corollary}

\textsc{Proof.}\ By Theorem \ref{thm:c4-4-7} $\Omega$ has a spherical harmonic development $\sum_{k=1}^{\infty} Y^{(k)}$ and $\Omega_0$ has the development $\sum_{k=1}^{\infty} Y_0^{(k)} = \sum_{k=1}^{\infty} \gamma_{k,0} Y^{(k)}$. Let $\Omega^{(m)} = \sum_{k=1}^{m} Y^{(k)}$ and $\Omega_0^{(m)} = \sum_{k=1}^{m} Y_0^{(k)}$. Then $\Omega_0^{(m)}$ is continuous and, by Theorem \ref{thm:c4-4-11},

\[
\begin{aligned}{\Omega_{0}(x^{\prime})-}&{{}\Omega_{0}^{(m)}(x^{\prime})=}\\ {}&{{}-\intop_{\Sigma_{n-1}}\left[\Omega(t^{\prime})-\Omega^{(m)}(t^{\prime})\right]\left[\frac{i\pi}{2}\text{sgn}{(x^{\prime}\cdot t^{\prime})}+\log{|x^{\prime}\cdot t^{\prime}|}\right]{d}t^{\prime}.}\\ \end{aligned}
\]

Thus, by Schwarz's inequality.

\[
\begin{aligned}&|\Omega_{0}(x^{\prime})-\Omega_{0}^{(m)}(x^{\prime})|\\
&\qquad\leq\left(\int_{\Sigma_{n-1}}|\Omega(t^{\prime})-\Omega^{(m)}(t^{\prime})|^{2}\,dt^{\prime}\right)^{1/2}\\
&\qquad\qquad\left(\int_{\Sigma_{n-1}}\left(\frac{\pi^{2}}{4}+|\log|x^{\prime}\cdot t^{\prime}||^{2}\right)dt^{\prime}\right)^{1/2}.\end{aligned}
\]

But the integral

\[
\int\limits_{\Sigma_{n-1}}\left(\frac{\pi^{2}}{4}\;+\;\left|\log|x^{\prime}\cdot t^{\prime}|\right|^{2}\right)d t^{\prime}
\]

is finite and is independent of $x'$, while $\int_{\Sigma_{n-1}}|\Omega(t')-\Omega^{(m)}(t')|^2dt' \to 0$ as $m \to \infty$. Consequently, $\{\Omega_0^{(m)}\}$ converges uniformly to $\Omega_0$. Since $\Omega_0^{(m)}$ is continuous, for $m = 1,2,3,\ldots$, it follows that $\Omega_0$ is continuous and the corollary is proved.

The principal value distributions $\Omega(x)/|x|^n$ just presented give an example of the distributions of type $(L^2, L^2)$ discussed in Chapter I, Section 3. As we shall see in Chapter VI, these distributions are also of type $(L^p, L^p)$ for $1 < p < \infty$.

Our next application concerns itself with yet another summability method for integrals. In the first chapter (see \eqref{eq:c1-1-12} we introduced the $\Phi$ means

\[
M_{\varepsilon,\Phi}(h)=\int\limits_{E_{n}}\Phi(\varepsilon x)h(x)d x
\]

of the integral $\int_{E_n} h$ whenever $\Phi$ belongs to $C_0$ and satisfies $\Phi(0) = 1$. When $\Phi(x) = e^{-|x|}$ we obtained the Abel means, while the Gauss-

Weierstrass means were obtained when $\Phi(x) = e^{-|x|^2}$. An equally important method of summability, Bochner--Riesz summability, is obtained when

\[
\Phi_{\delta}(x)=\Phi(x)=\left\{\begin{matrix}{(1-|x|^{2})^{\delta}}&{}&{\text{if }|x|\leq1}\\ {0}&{}&{\text{if }|x|>1}\\ \end{matrix}\right.
\]

for $\delta > 0$. In this case it is customary to put $\varepsilon = 1/R$ and, thus,

\[
{\cal M}_{\varepsilon,\Phi}(h)={\cal M}_{(1/R),\Phi}(h)=\int_{|x|\leq R}\left(1-\frac{|x|^{2}}{R^{2}}\right)^{\delta}h(x)\;d x.
\]

If $f$ belongs to $L^1(E_n)$, $\varphi = \hat{\Phi}$ and $h(x) = \hat{f}(x)e^{2\pi i t\cdot x}$ then (see Theorem \ref{thm:c1-1-16} of Chapter I)

\[
\begin{aligned}{S_{R}^{\delta}(t;f)}&{{}=S_{R}^{\delta}(t)=\int_{|x|\leq R}\hat{f}(x)e^{2\pi i x\cdot t}\bigg(1-\frac{|x|^{2}}{R^{2}}\bigg)^{\delta}d x}\\ {}&{{}=\int_{E_{n}}f(x)R^{n}\varphi(R(x-t))\;d x}\\ \end{aligned}
\]

are the Bochner--Riesz means of the integral $\int_{E_n} \hat{f}(x) e^{2\pi i x \cdot t} \, dx$ defining the inverse Fourier transform of $f$. If $\varphi$ belongs to $L^1(E_n)$ then, by Theorem \ref{thm:c1-1-18} of Chapter I, $S_R^\delta$ converges to $f$ in the norm, while if, in addition, $\varphi$ satisfies the hypotheses of Theorem \ref{thm:c1-1-25} of Chapter I we have

\[
\lim_{R\to\infty}S_{R}^{\delta}(t)=f(t)\quad\mathrm{a.e.}
\]

We shall show that this form of Fourier inversion does hold when $\delta > (n - 1)/2$. \footnote{The number $(n-1)/2$ is called the critical index for Bochner-Riesz\index{Bochner-Riesz summability@Bochner-Riesz summability} summability. We shall encounter it again in Chapter VII.\index{critical index@critical index}}

In order to do this we will calculate the Fourier transform $\varphi = \tilde{\Phi}$. The following identity involving Bessel functions is useful for performing this calculation:

\begin{lemma}
\label{lem:c4-4-13}
If $\mu > -\frac{1}{2}$ then
\[
J_{\mu+\nu+1}(t)=\frac{t^{\nu+1}}{2^{\nu}\Gamma(\nu+1)}\intop_{0}^{1}J_{\mu}(t s)s^{\mu+1}(1-s^{2})^{\nu}d s
\]

whenever $\nu > -1$ and t > 0.
\end{lemma}

\textsc{Proof.}\ If we develop $e^{its}$ into the power series $\sum_{j=0}^{\infty}(its)^{j}/j!$ it follows easily from the definition of $J_{k}(t)$, for real $k > -\frac{1}{2}$, that was given in the third section that

\begin{equation}\label{eq:c4-4-14}
J_{k}(t)=\sum_{j=0}^{\infty}(-1)^{j}\frac{(t/2)^{k+2j}}{j!\Gamma(j+k+1)}.\footnote{In deriving this formula it is useful to observe that the integral defining $J_{k}$ is equal to $2\int_{0}^{1}(\cos ts)(1-s^{2})^{(2k-1)/2}ds$. The coefficients of the power series \eqref{eq:c4-4-14} will then involve the expressions $\int_{0}^{1}s^{2j}(1-s^{2})^{(2k-1)/2}ds$ which can be easily evaluated by using the well-known relation $\Gamma(x)\Gamma(y)/\Gamma(x+y)=\int_{0}^{1}u^{x-1}(1-u)^{y-1}du$.}
\end{equation}

Thus, by \eqref{eq:c4-4-14},

\[
\begin{aligned}\int\limits_{0}^{1}J_{\mu}(ts)s^{\mu+1}(1-s^{2})^{\nu}ds&=\int\limits_{0}^{1}\biggl\{\sum_{j=0}^{\infty}(-1)^{j}\frac{(ts/2)^{\mu+2j}}{j!\;\Gamma(j+\mu+1)}\biggr\}s^{\mu+1}(1-s^{2})^{\nu}ds\\&=\sum_{j=0}^{\infty}(-1)^{j}\frac{(t/2)^{\mu+2j}}{j!\;\Gamma(j+\mu+1)}\frac{1}{2}\int\limits_{0}^{1}r^{\mu+j}(1-r)^{\nu}dr\\&=\frac{2^{\nu}\Gamma(\nu+1)}{t^{\nu+1}}\sum_{j=0}^{\infty}(-1)^{j}\frac{(t/2)^{\mu+\nu+1+2j}}{j!\;\Gamma(\mu+\nu+j+2)}\\&=\frac{2^{\nu}\;\Gamma(\nu+1)}{t^{\nu+1}}J_{\mu+\nu+1}(t)\end{aligned}
\]

and the lemma is proved.

\begin{theorem}
\label{thm:c4-4-15}
If $\Phi$ is the function defined by
\[
\Phi(t)=\left\{\begin{matrix}{(1-|t|^{2})^{\delta}}&{}&{\text{if }|t|\leq1}\\ {0}&{}&{\text{if }|t|>1,}\\ \end{matrix}\right.
\]

where $\delta>0$, then
\[
\begin{array}{r}{\hat{\Phi}(x)=\pi^{-\delta}\Gamma(\delta+1)|x|^{-(n/2)-\delta}J_{n/2+\delta}(2\pi|x|).}\end{array}
\]

\end{theorem}

\textsc{Proof.}\ By Theorem \ref{thm:c4-3-3}

\[
\begin{array}{r}{\hat{\Phi}(x)=2\pi|x|^{-[(n-2)/2]}\displaystyle\int_{0}^{1}(1-|s|^{2})^{\delta}J_{(n-2)/2}(2\pi|x|s)s^{n/2}\;d s.}\end{array}
\]

Now applying Lemma \ref{lem:c4-4-13} with $\nu = \delta$ and $\mu = (n - 2)/2$ we see that the last expression equals

\[
2\pi|x|^{-[(n-2)/2]}(2\pi)^{-\delta-1}2^{\delta}\Gamma(\delta+1)|x|^{-\delta-1}J_{(n-2)/2+\delta+1}(2\pi|x|),
\]

which clearly reduces to the desired expression for $\Phi(x)$.

It follows from \eqref{eq:c4-3-12} and \eqref{eq:c4-4-14} that $\varphi = \Phi$ belongs to $L^1(E_n)$ provided $\delta > (n - 1)/2$; moreover, $\varphi$ satisfies the additional hypotheses of Theorems \ref{thm:c1-1-18} and 1.25 of the first chapter and we obtain:

\begin{corollary}
\label{cor:c4-4-16}
If $f \in L^p(E_n)$, $1 \leq p < \infty$, and $S_R^\delta = \varphi_{(1/R)} * f$, where $\varphi_{(1/R)}(x) = R^n \varphi(Rx)$ for $R > 0$ and $\delta > (n - 1)/2$, then
(a) $\|S_R^\delta - f\|_p \to 0$ as $R \to \infty$;
(b) $\lim_{R \to \infty} S_R^\delta(x) = f(x)$ for all $x$ in the Lebesgue set of $f$ (in particular, this equality holds for almost every $x \in E_n$).
\end{corollary}

\section{Further Results}

\frp{5}{1} Theorem \ref{thm:c4-1-1} generalizes easily to a result describing the relation between the Fourier transform and an arbitrary nonsingular linear transformation of $E_n$ onto itself. Let $\sigma$ be such a transformation and define $\tilde{\sigma}$ to be the inverse of the adjoint of $\sigma$ (equivalently, $\tilde{\sigma}$ is the adjoint of the inverse of $\sigma$). The transformation $\tilde{\sigma}$ is called the contragredient of $\sigma$. If we define the action of $\sigma$ on the function $f$, $R_{\sigma}f$, by letting $(R_{\sigma}f)(x) = f(\sigma x)$ we then have the following result: If $f \in L^1(E_n)$ then $(R_{\sigma}f)^\wedge = |\det \sigma|^{-1} \times R_{\tilde{\sigma}} f$. That is, $\mathcal{F} R_{\sigma} = |\det \sigma|^{-1} R_{\tilde{\sigma}} \mathcal{F}$. The proof is no harder than that of Theorem \ref{thm:c4-1-1}: since the Jacobian in the change of variables $s = \sigma t$ is $|\det \sigma|^{-1}$, we have

\[
\begin{aligned}{(R_{\sigma}f)^{\wedge}(t)}&{{}=\intop_{E_{n}}e^{-2\pi i t\cdot x}f(\sigma x)\;{d}x\;=\;|\det\sigma|^{-1}\intop_{E_{n}}e^{-2\pi i(t\cdot\sigma^{-1}s)}f(s)\;{d}s}\\ {}&{{}=|\det\sigma|^{-1}\intop_{E_{n}}e^{-2\pi i(\tilde{\sigma}t\cdot s)}f(s)\;{d}s\;=\;|\det\sigma|^{-1}(R_{\tilde{\sigma}}\hat{f})(t).}\\ \end{aligned}
\]

\frp{5}{2} That orthogonal transformations and the Fourier transform are closely connected is further indicated by the following result (communicated to us by \textit{R}. Coifman): Suppose $\sigma$ is a continuous one-to-one transformation of $E_n$ onto itself such that $R_\sigma$ commutes with the Fourier transform; then $\sigma$ is an orthogonal transformation. To see this we first observe that, by assumption, $(R_\sigma \mathcal{F} f)(0) = (\mathcal{F} R_\sigma f)(0)$ for all $f \in L^1(E_n)$. Applying this equality to the function defined by $f(x) = g(x)e^{-2\pi i\sigma^{-1}(x) \cdot y}$ for $x \in E_n$, where $y \in E_n$ is fixed, we obtain

\[
0=\int_{E_{n}}[e^{-2\pi i x\cdot\sigma(0)}e^{-2\pi i\sigma^{-1}(x)\cdot y}g(x)\;-\;e^{-2\pi i x\cdot y}g(\sigma x)]\;d x.
\]

But, again by assumption, $\int_{E_n} e^{-2\pi i x \cdot y} g(\sigma x) \, dx = \int_{E_n} e^{-2\pi i x \cdot \sigma y} g(x) \, dx$; thus,

\[
0=\int\limits_{E_{n}}[e^{-2\pi i(x\cdot h+\sigma^{-1}(x)\cdot y)}-e^{-2\pi i x\cdot\sigma(y)}]g(x)d x
\]

for all $g$ in $L^{1}(E_{n})$, where $h = \sigma(0)$. Consequently, $e^{-2\pi i(x \cdot h + \sigma^{-1}(x) \cdot y)} = e^{-2\pi i x \cdot \sigma(y)}$ for all $x, y$ in $E_{n}$. Thus, $x \cdot h + \sigma^{-1}(x) \cdot y = x \cdot \sigma(y) + k$ for some integer $k$; however, the continuity of $\sigma$ implies $k$ is constant. Putting $x = 0$ we then have $\sigma^{-1}(0) \cdot y = k$ for all $y \in E_{n}$ which implies $k = 0$. Therefore, $h = \sigma(0) = 0$ since $\sigma^{-1}(0) = 0$. We have shown that $\sigma^{-1}(x) \cdot y = x \cdot \sigma(y)$ or, equivalently, $x \cdot y = \sigma(x) \cdot \sigma(y)$ for all $y \in E_{n}$. Hence, $\sigma$ is an isometry mapping 0 onto 0. But this clearly implies that $\sigma$ is an orthogonal transformation.

\frp{5}{3} Let $\beta(x, y)$, $x, y \in E_n$, be a nondegenerate real symmetric quadratic form (not necessarily positive). Then the mapping $t \to \hat{x}(t) = e^{2\pi i\beta(x, t)}$, where $x \in E_n$ is fixed, is a character\index{character@character} (by this we mean a continuous map of $E_n$ into the complex numbers of absolute value 1 such that $\hat{x}(s + t) = \hat{x}(s)\hat{x}(t)$ whenever $s$ and $t$ belong to $E_n$). With multiplication of all characters defined by letting $\hat{x}_1\hat{x}_2(t) = \hat{x}_1(t)\hat{x}_2(t)$ the collection of all characters is a group, the character group of $E_n$. The mapping $x \to \hat{x}$ is an isomorphism of $E_n$ into its character group. It is not hard to show that this mapping is onto; thus, we can ``identify'' $E_n$ with its character group. This identification clearly depends on the bilinear form $\beta$.

The definition of character and character group can obviously be extended to any locally compact abelian group $G$; however, it is not true in general that $G$ and its character group $\hat{G}$ are isomorphic. In abstract harmonic analysis the Fourier transform of a function $f \in L^{1}(G)$ is defined to be the function on $\hat{G}$ whose value at $\hat{x}$ is $\hat{f}(\hat{x}) = \int_{G} f(t) \overline{\hat{x}(t)} \, d\mu(t)$, where $\mu$ is Haar measure. In terms of the identification between $E_n$ and $\hat{E}_n$ described above we obtain a ``Fourier transform'' $T = T_\beta$ for each non-degenerate real symmetric quadratic form $\boldsymbol{\beta}: (T\hat{f})(x) = \int_{E_n} \overline{\hat{x}(t)} f(t) \, dt = \int_{E_n} e^{-2\pi i\beta(x,t)} f(t) \, dt$. When $\beta(x,t) = x \cdot t$ we obtain the Fourier transform we have been studying. In this chapter we exploited the close connection between the Fourier transform and the orthogonal group. In general, $T_\beta$ is similarly related to the orthogonal group associated with $\beta$: the group $0_\beta(n)$ of all linear transformations $\sigma$ of $E_n$ leaving the quadratic form $\beta$ invariant (that is, $\beta(\sigma x, \sigma y) = \beta(x, y)$ for all $x, y$ in $E_n$). If $\sigma \in 0_\beta(n)$ and $R_\sigma$ denotes the action of $\sigma$ on functions whose domain is $E_n$ (see 5.1) the argument used to establish Theorem \ref{thm:c4-1-1} also shows that $R_\sigma T_\beta = T_\beta R_\sigma$. Moreover, result 5.2 also extends to this situation: If $\sigma$ is a continuous one-to-one transformation of $E_n$ onto itself such that $R_\sigma$ commutes with $T_\beta$ then $\sigma \in 0_\beta(n)$.

\frp{5}{4} We have already stressed the significance of the group of dilations. The role of these transformations is further highlighted by considering the Mellin transform, in terms of which some of the results of this chapter can be reinterpreted.

Suppose $f$ is a complex-valued function on the interval $(0, \infty)$ which satisfies the integrability condition

\[
\int\limits_{0}^{\infty}\frac{\left|f(x)\right|}{x}dx<\infty.
\]

Then its Mellin\index{Mellin transform@Mellin transform} transform, $\mathcal{M}_{f}(s) = \mathcal{M}(s)$, is defined by

\[
\mathcal{M}(s)=\int\limits_{0}^{\infty}f(x)x^{s}\frac{dx}{x},
\]

for $\text{Re}(s) = 0$. If $f_1$ and $f_2$ both satisfy this integrability condition, then so does their multiplicative\index{convolution@convolution!multiplicative@multiplicative} convolution, $f_1 \cdot f_2$, defined by

\[
(f_{1}\cdot f_{2})(x)=\int\limits_{0}^{\infty}f_{1}(x/y)f_{2}(y)\frac{dy}{y},
\]

and we have $\mathcal{M}_{f_1 \cdot f_2}(s) = \mathcal{M}_{f_1}(s) \mathcal{M}_{f_2}(s)$. The Mellin transform is the multiplicative analogue of the Fourier transform, and many of its elementary properties can be deduced from the Fourier transform on $E_1$ by making the change of variables $x \to e^x$, taking $E_1$ to $(0, \infty)$. In this connection it may be of interest to state the following identities.

\begin{equation}
\tag{a}
\mathcal{M}_{f}(s)=2^{a+s-\mu-1}\Gamma\left(\frac{a+s}{2}\right)\bigg/\Gamma\left(\mu-\frac{a+s}{2}+1\right),
\end{equation}

if $f(x) = x^{a-\mu}J_\mu(x)$, with $a > 0$, $\mu > a - \frac{1}{2}$. The case when $\mu$ is integral or half-integral is essentially equivalent with Theorem \ref{thm:c4-4-1} (if we assume Theorem \ref{thm:c4-3-10} as known).

(b) Let $f_1(x) = x^{-\mu + 1} J_\mu(x)$, and $f_2(x) = x^{-2\mu - 2\nu + 1} (x^2 - 1)^\nu / 2^\nu \Gamma(\nu + 1)$, for $x \geq 1$, $f_2(x) = 0$, if $0 < x < 1$. Then $f(x) = (f_1 \cdot f_2)(x) = x^{-\mu - \nu} J_{\mu + \nu}(x)$, as long as $\nu > -1$, and $\mu > -\frac{1}{2}$. This identity is equivalent with the statement of Lemma \ref{lem:c4-4-13}. An alternative proof for that lemma can be carried out by showing that $\mathcal{M}_f(s) = \mathcal{M}_{f_1}(s) \mathcal{M}_{f_2}(s)$.

For a proof of (a) and further facts concerning the Mellin transform, see Titchmarsh [2].

\frp{5}{5} In the second chapter we showed that a harmonic function has continuous derivatives of all orders (see Theorem \ref{thm:c2-1-7} in that chapter). It is an immediate consequence of the Poisson integral representation of harmonic functions (see Theorem \ref{thm:c2-1-10} and Corollary \ref{cor:c2-1-11}) together with Theorem \ref{thm:c4-2-10} that a harmonic function is real-analytic.

\frp{5}{6} Lemma \ref{lem:c4-2-11} is equivalent to the statement that the only differential polynomials with constant coefficients that commute with rotations are

polynomials in the Laplacian. More precisely, the lemma is equivalent to the following statement, when n > 1.

If $P(D)$, $D = (\partial/\partial x_1, \ldots, \partial/\partial x_n)$, is a differential polynomial on $E_n$ with constant coefficients, then $P(D)R_\sigma = R_\sigma P(D)$ for all rotations $\sigma$ if and only if $P(D) = c_0I + c_1\Delta + \cdots + c_m\Delta^m$, where $\Delta = \partial^2/\partial x_1^2 + \cdots + \partial^2/\partial x_n^2$ is the Laplace operator.

\frp{5}{7} We defined the ultraspherical polynomials $P_k^\lambda$ in terms of the generating function $(1 - 2rt + r^2)^{-\lambda}$. When $\lambda = 0$, however, it is more appropriate to consider the Tchebichef\index{Tchebichef polynomials@Tchebichef polynomials} polynomials $T_k$ whose generating function is

\[
\log\left(1-2r t+r^{2}\right)=-2\sum_{k=1}^{\infty}k^{-1}T_{k}(t)r^{k}.
\]

If this is done it is not necessary to exclude the two-dimensional case, as was done often in the second section. For example, Theorem \ref{thm:c4-2-14} is valid when n = 2 if we employ the Tchebichef polynomials.

\frp{5}{8} An important connection between the group $\mathrm{SO}(n)$ of rotations on $E_n$ and spherical harmonics arises in the study of the unitary representations of this group. The mapping $\sigma \to R_{\sigma^{-1}}$, where $R_{\sigma^{-1}}$ is the operator on $L^2(\Sigma_{n-1})$ defined by $(R_{\sigma^{-1}}f)(x) = f(\sigma^{-1}x)$, is a unitary representation of $\mathrm{SO}(n)$. When we restrict $R_{\sigma^{-1}}$ to the subspace $\mathcal{H}_k$, $k = 0, 1, 2, \ldots$, we obtain an irreducible representation of $\mathrm{SO}(n)$. When the dimension $n$ is 3 all the irreducible representations of $\mathrm{SO}(n)$ can be obtained in this way. The group $\mathrm{SO}(n-1)$ can be identified in a natural way with the subgroup $G(n, x')$ of $\mathrm{SO}(n)$ consisting of those rotations $\sigma$ leaving a fixed point $x' \in \Sigma_{n-1}$ fixed. We see from Lemma \ref{lem:c4-2-8} that in each of the spaces $\mathcal{H}_k$ there exists a nonzero vector (viz., $Z_x'$) left invariant by all the operators $R_\sigma$, where $\sigma$ ranges over the subgroup $G(n, x')(\sim \mathrm{SO}(n-1))$. Among all the irreducible representations of $\mathrm{SO}(n)$, the ones described above are characterized, up to a unitary equivalence, by the existence of such a vector that is invariant under the action of $\mathrm{SO}(n-1)$ (see Weyl [1] and Boerner [1]).

\frp{5}{9} Corollary \ref{cor:c4-4-12} extends to the case $\Omega \in L^p(\Sigma_{n-1})$, $p > 1$. The proof is not essentially different: instead of Schwarz's inequality one can employ Hölder's inequality and, instead of using the $L^2$-norm for approximating functions by spherical harmonics, one can avail himself of the $L^p$ norm. A further refinement of this argument shows that the result still holds when $|\Omega| \log^+ |\Omega|$ is integrable (where $\log^+ x = \log x$ if $x \geq 1$, while $\log^+ x = 0$ if $0 \prec x < 1$).

\frp{5}{10} The analogues of the subspaces $\mathfrak{G}_k$ can obviously be defined for $L^p(E_n), 1 \leq p \leq \infty$. Let us remark that if $f_0(|x|)P(x)$ belongs to $L^p(E_n)$,

1 ≤ p ≤ 2, then its Fourier transform $\hat{f}$ is given by the formula \eqref{thm:c4-3-10}, where the convergence is in the $L^p(E_n)$ norm (see \eqref{eq:c5-1-2} of Chapter V). The estimates in Lemma \ref{lem:c4-3-11} for the Bessel functions then have the following consequences: (a) $\hat{f}(x)$ is equivalent to a function continuous when $x \neq 0$, if $1 \leq p < 2n/(n+1)$; (b) when $k < (1/p - \frac{1}{2})n - \frac{1}{2}$, then $\hat{f}$ is of class $C^{(k)}$ away from the origin; (c) when p = 1, the condition can be modified to require $k \leq (n-1)/2$ (results of this type were communicated to us by N. Varapoulous). These facts are simple corollaries of Lemma \ref{lem:c4-3-11}, Hölder's inequality and the relation $(t^{-k}J_k(t))' = -t^{-k}J_{k+1}(t)$ (see \eqref{eq:c4-3-2}).

\bibnoteshead

We introduced only those properties of Bessel functions that were needed for our study of the Fourier transform. For a thorough treatment of this subject see Watson [1] or the Bateman Manuscript Project [1], Vol. 2. We also refer the reader to the latter for further details concerning spherical harmonics, Gegenbauer and Tchebichet polynomials. For earlier studies of the action of the Fourier transform on the spaces $\mathbb{Q}_k$ see Hecke [1], Bochner [6], and Herz [1]. For another proof of Theorem \ref{thm:c4-4-1} see Calderón [3]. The Fourier transforms of the kernels $K(x) = \Omega(x')/|x|^n$ were first thoroughly studied by Mihlin [1], and Calderón and Zygmund [5]. The properties of the rotation group can be exploited further to obtain various results related to those in this chapter. See Vilenkin [1], Coifman and G. Weiss [1].

